% ---------------------------------------------
% Copyright Ben Paul Wise. All rights reserved.
% ---------------------------------------------
% process with TeXworks as pdfLaTeX
%
% Preliminary problem statement for the acv hierarchical Bayesian work.
% Algebra verified in maxima/acv_prelim.mac

\documentclass[english, 11pt]{report}

% ---------------------------------------------
% Copyright Ben Paul Wise. All rights reserved.
% ---------------------------------------------
% process with TeXworks as pdfLaTeX+MakeIndex+BibTex


% This file is \input by main.tex as the STYLE MODULE, so the class line is
% commented out here and the \documentclass (report, for chapters) lives in
% main.tex; everything below applies unchanged.
%\documentclass[english, 10pt]{article}
% `english' is necessary for babel, not `letter' or `a4paper'
% default is 10pt unless specified otherwise.
% I reserve 12 pt for legible final documents


%\usepackage[a4paper]{geometry}
\usepackage[letterpaper]{geometry}

\usepackage{geometry}
\geometry{verbose,lmargin=2.0cm,rmargin=2.0cm}

\setlength{\parindent}{2em}
\setlength{\parskip}{0.75em}

\usepackage{fancyhdr}
\pagestyle{fancy}
\usepackage{color}
\definecolor{mypurple}{RGB}{128, 0, 128}
\definecolor{myTeal}{RGB}{0, 128, 128}
\definecolor{myDarkGreen}{RGB}{0, 92, 0}
\definecolor{myLightGreen}{RGB}{0, 192, 0}
\definecolor{myDarkRed}{RGB}{192, 0, 0}
\definecolor{myDarkBlue}{RGB}{0, 0, 192}
\definecolor{myBeige}{RGB}{255, 255, 192}

\definecolor{GroupWBlue1}{RGB}{75, 99, 175} 
\definecolor{GroupWGrey1}{RGB}{145, 147, 150} 

\definecolor{schrift}{RGB}{0,73,114}

\usepackage{amsmath} %% get AMS standards for theorems, e.g. the hollow QED box,  \qedsymbol
%% But I like the solid "tombstone", Halmos, whatever:
\newcommand{\qedsolid}{\rule{1ex}{1.62ex}}
\usepackage{amssymb} % \sum


\usepackage{empheq} % colored equation-boxes
\usepackage[most]{tcolorbox}
\tcbset{colback=white!90!yellow, % i.e. 90% and 10% yellow
        colframe=green!36!blue, % i.e. 64% green and 36% blue, `schrift'
        highlight math style= {enhanced, %<-- needed for the ’remember’ options
            colframe=red,colback=red!10!white,boxsep=0pt}
        }
		
		
% setup colored and boxed quotes	
% this would give a yellowish background inside a reddish frame:
% colback=yellow!10!white, colframe=red!50!black, 

%\usepackage[most]{tcolorbox}     
\newtcolorbox{myquote}{
colback=yellow!10!white,
colframe=blue!50!black,
grow to right by=-20mm,
grow to left by=-30mm,
%boxrule=2pt,
boxsep=0pt,
breakable} \makeatletter

\newcommand{\blockquote}[1]{  \begin{myquote}  #1  \end{myquote}  }


% These are usually 'schrift' not 'GroupWBlue1'
\usepackage{hyperref}   % clickable links into or outfrom this text
\hypersetup{
    colorlinks=true,
    linkcolor=GroupWBlue1,
    filecolor=magenta,      
    urlcolor=cyan,
    citecolor=myTeal
}

\numberwithin{equation}{section} % number equations as (sec.eq), not just (eq)

% Good reference on tables in LaTeX:
% https://en.wikibooks.org/wiki/LaTeX/Tables 

\usepackage{graphicx}
% parameters for the 'figure box'
\fboxsep=1mm %padding thickness
\fboxrule=2pt %border thickness

\usepackage{sectsty}

% These are usually 'schrift' not 'GroupWBlue1' or 'GroupWGrey1'
\sectionfont{\color{GroupWBlue1}}  % sets colour of sections
\subsectionfont{\color{GroupWBlue1}}
\subsubsectionfont{\color{GroupWBlue1}}

\usepackage{lastpage}

% The following two are needed to shade individual cells,
% without the `feature' of over-writing the cell boundaries.
% However, they seem to clash with empheq, so they are commented out.
% this means that \cellcolor{color}{stuff} cannot be used; see revisions just prior to 2022-03-27 for colored cells.
%\usepackage[table]{xcolor}% http://ctan.org/pkg/xcolor, https://tex.stackexchange.com/questions/50349/color-only-a-cell-of-a-table
%\usepackage{hhline}% http://ctan.org/pkg/hhline, https://tex.stackexchange.com/questions/25062/cell-shading-without-removing-lines


\usepackage[T1]{fontenc} % necessary for babel
\usepackage{babel}
\usepackage{url} % so bitex can handle URL


%\rhead{Basic CP}   % default rhead is section
%\lhead{B. P. Wise} % default lhead is subsection

% Some people say that ragged right edges make it easier to read.
\raggedright

% ---------------------------------------------
% Copyright Ben Paul Wise. All rights reserved.
% ---------------------------------------------



% The report class numbers sections as chapter.section.  This document uses
% no chapters, so the chapter counter stays at zero and the sections come
% out as 0.1, 0.2 and so on.  Numbering them by section alone gives 1, 2, 3.
% Equations follow, since prolog.tex sets \numberwithin{equation}{section}.
\renewcommand{\thesection}{\arabic{section}}

\def\myversion{v01a}

%\rhead{Basic VI-NCP} % default is section
%\lhead{B. P. Wise} % default is subsection

\lfoot{\color{myDarkRed}\bf{DRAFT \myversion}}
\cfoot{{\thepage} of {\pageref*{LastPage}}} % asterix suppressed hyperlink
%\rfoot{\today}
\rfoot{\copyright \  Ben P. Wise}

\begin{document}

\begin{titlepage}
\title{Recovering a Sparse Matrix from Paired Observations\\
       A Preliminary Problem Statement}
\author{Ben Paul Wise}
\date{\today}

\maketitle

\begin{abstract}
summary goes here
\end{abstract}

\end{titlepage}


\clearpage
\tableofcontents
\clearpage
\listoffigures
\clearpage
\listoftables
\clearpage


\raggedright

\section{Purpose}

This document states one small problem completely, and solves it as far as algebra
will carry it. The problem is the first step toward a larger programme: inference of
vector time-series models in which the structure of the model is determined by the
data rather than assumed beforehand. Nothing about that larger programme appears
here. Later documents will add one element at a time.

The problem treated here is the following. We observe pairs of vectors, an input and
an output, and we believe the output is a fixed linear function of the input, plus
measurement error. We want the matrix of that linear function. We expect most of its
entries to be zero, and we do not know beforehand which ones.

\section{Bayes's theorem and the estimate we shall use}

Let $A$ denote a proposition and $B$ a body of evidence. Bayes's theorem starts
with the basic definition of conditional probability

\begin{equation}
\tcboxmath{ \begin{array}[t]{c}
P(A \mid B)  P(B) \;=\;  P(A,B)  \;=\; P(B \mid A)\, P(A)
\end{array}}
\end{equation}

and derives the formula for updating beliefs given new evidence: 

\begin{eqnarray}
P(A \mid B) &=& \frac{P(B \mid A)\, P(A)}{P(B)}
\label{eq:bayes_A}
\\
P(B) &=& \int P(B \mid a)\, P(a)\, da
\label{eq:bayes_B}
\end{eqnarray}



where the integral runs over every proposition $a$ that might have produced the
evidence, and becomes a sum when that set is countable.

The denominator $P(B)$ does not involve $A$. It is the same number for every
candidate $A$ under consideration. Consequently, if our purpose is to find the
candidate $A$ that makes $P(A \mid B)$ as large as possible, we may work with the
numerator alone and never evaluate $P(B)$ at all:

\begin{equation}
\arg\max_{A} P(A \mid B) \;=\; \arg\max_{A} \; P(B \mid A)\, P(A)
\label{eq:mple}
\end{equation}

This is maximum posterior likelihood estimation, and it is what this document does.
The saving is not incidental. The integral defining $P(B)$ runs over every matrix
that might have produced the data, and we have no way to evaluate it.

Two warnings apply to \eqref{eq:mple} and are taken up in
Sections~\ref{sec:magnitudes} and~\ref{sec:noise}. First, the largest value of a
probability density is not always attained; a density can increase without bound.
Second, the position of the largest value depends on how the parameters are written,
since a density is defined relative to a measure. Both matter here, and both are
resolved by stating what we know rather than by mathematical device.

\section{Notation}

\begin{center}
\begin{tabular}{ll}
$N$ & the dimension of the vectors; $A$ is $N \times N$ \\
$T$ & the number of observation pairs \\
$n$ & $NT$, the number of scalar measurements \\
$q$ & $N^2$, the number of cells in $A$ \\
$X_{jt}$ & component $j$ of the input vector at observation $t$ \\
$Y_{it}$ & component $i$ of the output vector at observation $t$ \\
$Z_{it}$ & the value of $Y_{it}$ the matrix predicts \\
$\mathcal{S}$ & the set of cells $(i,j)$ at which $A_{ij} \neq 0$ \\
$M$ & the number of cells in $\mathcal{S}$ \\
$p$ & the probability that any one cell is non-zero \\
$\sigma$ & the standard deviation of the measurement error \\
$R(A)$ & the residual sum of squares \\
\end{tabular}
\end{center}

\section{The space of matrices}
\label{sec:matrixspace}

The unknown $A$ is a square array of $N$ rows and $N$ columns. Each cell holds a real
number, positive or negative, or else holds nothing at all, which we write as zero.
A candidate is therefore two things taken together: a choice of which cells are
occupied, and a choice of value for each occupied cell.

The first choice is discrete. There are $2^{q}$ possible patterns, where $q = N^2$.
The second is continuous, and once the pattern is fixed it involves $M$ real numbers.
Any prior over $A$ must give both.

\subsection{The pattern of occupied cells}

We suppose that each of the $q$ cells is occupied or not independently of the others,
with the same probability $p$ of being occupied. A pattern with $M$ occupied cells
then has probability

\begin{equation}
P(\mathcal{S} \mid p) \;=\; p^{M} \, (1-p)^{\,q-M}
\label{eq:patternprior}
\end{equation}

The second factor is necessary. Without it the expression does not sum to unity over
the $2^{q}$ patterns, and the comparison between a sparse matrix and a dense one is
not a comparison between probabilities.

We do not know $p$. The hyper-prior over it is taken to be a Beta distribution with
parameters $\alpha$ and $\beta$,

\begin{equation}
P(p) \;=\; \frac{p^{\alpha-1} (1-p)^{\beta-1}}{B(\alpha,\beta)}
\label{eq:pprior}
\end{equation}

where $0 \le p \le 1$. The mean of this distribution is
$\alpha/(\alpha+\beta)$. Setting $\alpha = \beta = 1$ gives the uniform
distribution on $[0,1]$, which is the statement that nothing is known about $p$
beyond its range. Choosing $\alpha/(\alpha+\beta)$ below one half states an
expectation that the matrix is sparse without excluding a dense one, and the sum
$\alpha + \beta$ governs how firmly that expectation is held. The general form is
carried through the algebra below; Section~\ref{sec:patterncost} examines what
influence it can have.

All algebra and integrals in this paper were derived and verified using the open-source symbolic
mathematics package Maxima version 5.49.0.

\subsection{The magnitudes of the occupied cells}
\label{sec:magnitudes}

For the value in an occupied cell we should like a statement that does not depend on
the units in which $X$ and $Y$ are measured. The usual such statement is that the
magnitude is equally likely to fall in any interval of fixed width on a logarithmic
scale: as likely between $0.01$ and $0.1$ as between $10$ and $100$. Written as a
density in $a$, that statement is proportional to $1/|a|$.

Taken over the whole real line, $1/|a|$ is not a probability density. Its integral
diverges at zero and again at infinity,

\begin{eqnarray}
\int_{\varepsilon}^{1} \frac{da}{a} &=& -\log \varepsilon
\;\xrightarrow[\varepsilon \to 0]{}\; \infty
\label{eq:diverge_A}
\\ \nonumber
\\
\int_{1}^{K} \frac{da}{a} &=& \log K
\;\xrightarrow[K \to \infty]{}\; \infty
\label{eq:diverge_B}
\end{eqnarray}

so it has no normalising constant. Densities of this kind are used, and are
sometimes harmless, but here they fail in two separate ways.

\paragraph{The posterior would have no maximum.}
Hold the pattern fixed. As will be shown in Section~\ref{sec:criterion}, the quantity
to be minimised contains the term $\sum \log|A_{ij}|$, summed over occupied cells.
Let one coefficient approach zero while the others are held fixed. The residual sum
of squares is a continuous quadratic function of the coefficients, so it approaches
the finite value it takes when that coefficient equals zero. The term $\log|A_{ij}|$,
by contrast, decreases without bound. The quantity to be minimised therefore has no
lower bound, and no maximising matrix exists.

The infimum is approached as the coefficient tends to zero, and at zero the candidate
is no longer a member of the set being searched: a coefficient equal to zero means
the cell is unoccupied, which is a different pattern with $M$ smaller by one. More
data does not repair this. Increasing $n$ enlarges the term measuring fit relative to
the others, but $\log|A_{ij}|$ still diverges.

\paragraph{Patterns of different size could not be compared.}
Suppose we avoided the first difficulty by integrating over the coefficient values
instead of maximising over them. A density without a normalising constant is
determined only up to an arbitrary positive factor $c$: the density $c/|a|$ expresses
the same ignorance for every $c$. A pattern with $M$ occupied cells then contributes
$c^{M}$, and the comparison of a pattern of size $M$ with one of size $M+1$ is
multiplied by $c$. Since $c$ has no determined value, any conclusion whatever about
the sparsity of $A$ can be produced by choosing $c$ suitably.

The general rule is this. A density without a normalising constant may be used when
the number of factors it contributes is the same for every candidate, since the
unknown constant then cancels from every comparison. For the coefficients of $A$ that
number is $M$, and $M$ varies from candidate to candidate.

\paragraph{The remedy.}
We keep the logarithmic statement but confine it to a stated range of magnitudes.
The analyst using this method knows, for the problem at hand, the smallest and the
largest coefficient magnitudes worth entertaining. Call them $a_{\min}$ and
$a_{\max}$, and write

\begin{equation}
L_{A} \;=\; \log \frac{a_{\max}}{a_{\min}}
\label{eq:LA}
\end{equation}

The prior for one occupied cell is then

\begin{equation}
P(A_{ij}) \;=\; \frac{1}{2 L_{A} \, |A_{ij}|}
\label{eq:coefprior}
\end{equation}

where $a_{\min} \le |A_{ij}| \le a_{\max}$, the factor two accounting for the two
signs. This integrates to unity. Both
difficulties disappear: $\log|A_{ij}|$ is bounded below by $\log a_{\min}$, and the
normalising constant is determined, so patterns of different size may be compared.

The invariance that recommended $1/|a|$ in the first place survives. A change in the
units of the coefficients multiplies $a_{\min}$ and $a_{\max}$ by the same factor and
leaves $L_{A}$ unchanged. The quantity $L_{A}$ has a plain reading: it is the width of
the admissible range of magnitudes, measured in natural logarithms. A range of ten
powers of ten gives $L_{A} \approx 23$.

Collecting \eqref{eq:patternprior} and \eqref{eq:coefprior}, the prior over the whole
matrix, given $p$, is

\begin{equation}
P(A \mid p) \;=\; p^{M} (1-p)^{\,q-M} \, (2 L_{A})^{-M}
\prod_{(i,j) \in \mathcal{S}} \frac{1}{|A_{ij}|}
\label{eq:Aprior}
\end{equation}

\section{The data and the measurement error}
\label{sec:noise}

The data $B$ consist of $T$ pairs of $N$-dimensional real vectors. The inputs $X$ are
taken to be exact. All error is in the outputs. The matrix predicts

\begin{equation}
Z_{it} \;=\; \sum_{j=1}^{N} A_{ij} X_{jt}
\label{eq:Z}
\end{equation}

and we suppose the discrepancies $Y_{it} - Z_{it}$ to be independent of one another,
each drawn from a normal distribution of mean zero and standard deviation $\sigma$.
Then

\begin{equation}
P(B \mid A, \sigma) \;=\; \prod_{i=1}^{N} \prod_{t=1}^{T}
\frac{1}{\sigma \sqrt{2\pi}}
\exp\!\left( -\frac{(Z_{it} - Y_{it})^{2}}{2\sigma^{2}} \right)
\;=\; (2\pi)^{-n/2} \, \sigma^{-n} \exp\!\left( -\frac{R(A)}{2\sigma^{2}} \right)
\label{eq:likelihood}
\end{equation}

where

\begin{equation}
R(A) \;=\; \sum_{i=1}^{N} \sum_{t=1}^{T} \left( Y_{it} - Z_{it} \right)^{2}
\label{eq:R}
\end{equation}

and $n = NT$. The factors of $\sqrt{2\pi}$ contribute $(2\pi)^{-n/2}$, which does not
involve $A$, $p$ or $\sigma$. It is the same number for every candidate and may be
dropped throughout without affecting which candidate is best.

\subsection{The scale of the error}

We do not know $\sigma$. The natural statement of ignorance about a scale is again
the logarithmic one, proportional to $1/\sigma$, and again it has no total mass.

Here the consequence differs from the one in Section~\ref{sec:magnitudes}. Every
candidate matrix carries exactly one factor of the prior for $\sigma$, so the
arbitrary constant does cancel. The difficulty lies elsewhere, and it is clearest in
the integrand itself. Combining the prior $1/\sigma$ with the likelihood
\eqref{eq:likelihood}, and setting aside factors that do not involve $\sigma$, the
integral to be performed is

\begin{equation}
\int_{0}^{\infty} \sigma^{-(n+1)}
\exp\!\left( -\frac{R}{2\sigma^{2}} \right) d\sigma
\label{eq:sigmanaive}
\end{equation}

The factor $\sigma^{-(n+1)}$ is by itself not integrable at the origin. What makes
\eqref{eq:sigmanaive} converge there is the exponential, which drives the integrand
to zero as $\sigma$ decreases. The reach of that suppression is governed by $R$. The
integrand rises to a single maximum at

\begin{equation}
\sigma^{2} \;=\; \frac{R}{n+1}
\label{eq:peak}
\end{equation}

and falls away on both sides. As $R$ decreases, the maximum moves toward the origin
and its height, proportional to $R^{-(n+1)/2}$, grows without limit. The suppression
that made the integral finite is confined to an ever smaller neighbourhood of zero,
exposing more and more of the singularity in $\sigma^{-(n+1)}$. Performing the
integral,

\begin{equation}
\int_{0}^{\infty} \sigma^{-(n+1)}
\exp\!\left( -\frac{R}{2\sigma^{2}} \right) d\sigma
\;=\; \tfrac{1}{2}\, \Gamma\!\left( \tfrac{n}{2} \right)
\left( \frac{2}{R} \right)^{n/2}
\label{eq:sigmaunbounded}
\end{equation}

which increases without bound as $R$ approaches zero. The divergence occurs at
$\sigma = 0$, and it appears in the answer as the factor $R^{-n/2}$. Stated in words:
as the fit improves, an ever wider range of very small error scales becomes
consistent with the data, and the accumulated weight of those scales has no upper
limit.

The divergence is not a defect of the algebra. It is the formulation reporting that
nothing has been said about how small the error can be. With $1/\sigma$ taken over
$(0,\infty)$, a residual of $10^{-12}$ is admitted as evidence of an instrument
accurate to $10^{-12}$, and no part of the model contradicts that reading. Stating
$s_{0}$ contradicts it.

Whenever the data can be fitted exactly, therefore, the exact fit dominates every
other candidate however large its pattern.

Whether an exact fit is available depends on $N$ and $T$. Row $i$ of $A$ is
constrained only by row $i$ of $Y$, which supplies $T$ equations, and it contains as
many free coefficients as that row has occupied cells. Two regimes are worth
recording:

\begin{center}
\begin{tabular}{lccl}
 & $N=10,\ T=20$ & $N=200,\ T=10$ & \\
\hline
cells in $A$      & $100$  & $40{,}000$ & \\
measurements $n$  & $200$  & $2{,}000$  & \\
equations per row & $20$   & $10$       & \\
coefficients available per row & $10$ & $200$ & \\
exact fit possible & no & yes, easily & \\
\end{tabular}
\end{center}

In the first regime every row is over-determined and no exact fit exists. In the
second, a pattern with ten occupied cells in each row reproduces the data exactly,
and there are a great many such patterns. The second regime is the one we must be
able to handle, so the difficulty cannot be set aside.

The remedy is the one already applied to the coefficients. The analyst who can state
the range of credible coefficient magnitudes can equally state the range of credible
noise levels: $s_{0}$, below which the instrument cannot measure, and $s_{1}$, above
which the data would be worthless. Writing

\begin{eqnarray}
L_{\sigma} &=& \log \frac{s_{1}}{s_{0}}
\label{eq:sigmaprior_A}
\\
P(\sigma) &=& \frac{1}{L_{\sigma}\, \sigma}
\label{eq:sigmaprior_B}
\\
s_{0} &\le& \sigma \le s_{1}
\label{eq:sigmaprior_C}
\end{eqnarray}



gives a proper density, and, as Section~\ref{sec:sigmaint} shows, a finite result at
$R = 0$. The statement $\sigma \ge s_{0}$ is a statement about the apparatus, not a
mathematical convenience: it says that the measurements are not exact, and how
inexact they are known to be.

\section{The posterior for $A$, $p$ and $\sigma$}

Multiplying the prior \eqref{eq:Aprior}, the hyper-priors on $p$ and $\sigma$, and
the likelihood \eqref{eq:likelihood}, and discarding the constant $(2\pi)^{-n/2}$,
the posterior for the three unknowns together is proportional to

\begin{equation}
P(A, p, \sigma \mid B) \;\propto\;
p^{M} (1-p)^{\,q-M}
\cdot \frac{(2L_{A})^{-M}}{\prod_{\mathcal{S}} |A_{ij}|}
\cdot \frac{1}{L_{\sigma}\,\sigma}
\cdot \sigma^{-n} \exp\!\left( -\frac{R(A)}{2\sigma^{2}} \right)
\label{eq:joint}
\end{equation}

The quantities $p$ and $\sigma$ are not of interest in themselves. We remove them by
integration, which is the step that makes the estimate of $A$ depend on the data
alone and not on values we should otherwise have to guess.

\subsection{Removing $p$}

Only the first factor of \eqref{eq:joint} involves $p$. With the hyper-prior
\eqref{eq:pprior},

\begin{equation}
\frac{1}{B(\alpha,\beta)} \int_{0}^{1}
p^{\,M+\alpha-1} (1-p)^{\,q-M+\beta-1} \, dp
\;=\; \frac{B(M+\alpha,\; q-M+\beta)}{B(\alpha,\beta)}
\label{eq:pint}
\end{equation}

At $\alpha = \beta = 1$ this reduces to

\begin{equation}
B(M+1,\; q-M+1) \;=\; \frac{M!\,(q-M)!}{(q+1)!}
\;=\; \frac{1}{(q+1) \binom{q}{M}}
\label{eq:pintuniform}
\end{equation}

whose negative logarithm is $\log \binom{q}{M} + \log(q+1)$. The first term is the
number of natural units of information required to say which $M$ of the $q$ cells are
occupied. This is Occam's razor in a form that can be computed: a pattern must reduce
the residual by enough to account for the information needed to describe it.

In either case the result depends on the pattern only through its size. Write

\begin{equation}
\Pi(M) \;=\; -\log \frac{B(M+\alpha,\; q-M+\beta)}{B(\alpha,\beta)}
\label{eq:Pi}
\end{equation}

for the contribution of the pattern to the quantity to be minimised.

\subsection{How much the pattern term can do}
\label{sec:patterncost}

Since $\Pi$ depends only on $M$, its whole effect is contained in the amount by which
it grows when one more cell is occupied. That increment has a simple closed form. The
ratio of successive values of the Beta function is

\begin{equation}
\frac{B(M+\alpha,\; q-M+\beta)}{B(M+1+\alpha,\; q-M-1+\beta)}
\;=\; \frac{q-M-1+\beta}{M+\alpha}
\label{eq:ratio}
\end{equation}

so that

\begin{equation}
\Pi(M+1) - \Pi(M) \;=\; \log \frac{q-M-1+\beta}{M+\alpha}
\label{eq:increment}
\end{equation}

Two consequences follow, and both matter for what a Beta hyper-prior can accomplish.

First, the increment depends on $\beta$ only through its logarithm. To raise the
charge for one more occupied cell by $D$ natural units above its value under the
uniform hyper-prior requires $\beta$ to grow by a factor of about $e^{D}$. A
hyper-prior that is merely sceptical of dense matrices therefore changes the
increment very little. With $q = 81$ and $M = 27$, the increment is $0.657$ at
$\alpha = \beta = 1$, and $1.281$ at $\alpha = 1,\ \beta = 50$, a hyper-prior whose
mean is $1/51$. Supplying a further six natural units would require $\beta$ of order
$10^{4}$, which asserts that the expected number of occupied cells in the whole
matrix is less than one hundredth.

Second, the increment falls as $M$ rises, so the pattern term resists the first few
occupied cells more than the last. Its total contribution across the whole matrix is
bounded by $q \log q$, of order four hundred natural units for $q = 81$, whereas the
term measuring fit varies over hundreds of units for a single pattern. The
hyper-prior on $p$ is therefore a refinement, not a control. Anything that must be
prevented has to be prevented elsewhere.

The posterior mean of $p$, should it be wanted, is $(M+\alpha)/(q+\alpha+\beta)$.

\subsection{Removing $\sigma$}
\label{sec:sigmaint}

The remaining factors involving $\sigma$ give

\begin{equation}
I(R) \;=\; \frac{1}{L_{\sigma}}
\int_{s_{0}}^{s_{1}} \sigma^{-(n+1)}
\exp\!\left( -\frac{R}{2\sigma^{2}} \right) d\sigma
\label{eq:sigmaintdef}
\end{equation}

The substitution $u = 1/\sigma^{2}$ turns this into an incomplete gamma integral.
With $\Gamma(m,x) = \int_{x}^{\infty} u^{m-1} e^{-u}\, du$ and $m = n/2$,

\begin{equation}
I(R) \;=\; \frac{1}{2 L_{\sigma}}
\left( \frac{2}{R} \right)^{n/2}
\left[ \Gamma\!\left( \frac{n}{2}, \frac{R}{2 s_{1}^{2}} \right)
     - \Gamma\!\left( \frac{n}{2}, \frac{R}{2 s_{0}^{2}} \right) \right]
\label{eq:sigmaint}
\end{equation}

Two limits show what the bounds accomplish. If $R$ is comfortably larger than
$2 s_{0}^{2}$, the second incomplete gamma function is negligible and the first is
close to $\Gamma(n/2)$, so

\begin{equation}
I(R) \;\approx\; \frac{\Gamma(n/2)}{2 L_{\sigma}}
\left( \frac{2}{R} \right)^{n/2}
\label{eq:sigmaordinary}
\end{equation}

which is \eqref{eq:sigmaunbounded}, the result obtained from the unbounded prior
divided by $L_{\sigma}$. The bounds cost nothing when
the fit is ordinary. If instead $R$ approaches zero, the factor $R^{-n/2}$ is
cancelled exactly by the behaviour of the bracket, and

\begin{equation}
\lim_{R \to 0} I(R) \;=\; \frac{1}{n L_{\sigma}}
\left( s_{0}^{-n} - s_{1}^{-n} \right)
\label{eq:sigmalimit}
\end{equation}

a finite number. A perfect fit is now merely very good, not infinitely good, and it
must still account for the size of its pattern.

\section{The criterion for the matrix}
\label{sec:criterion}

Collecting \eqref{eq:joint}, \eqref{eq:pint} and \eqref{eq:sigmaint}, the posterior
for the matrix alone is proportional to

\begin{equation}
P(A \mid B) \;\propto\;
e^{-\Pi(M)}
\cdot \frac{(2L_{A})^{-M}}{\prod_{\mathcal{S}} |A_{ij}|}
\cdot I\!\left( R(A) \right)
\label{eq:posteriorA}
\end{equation}

It is convenient to work with the negative logarithm, which is to be made as small as
possible. Discarding terms that do not involve $A$,

\begin{equation}
\tcboxmath{ \begin{array}[t]{c}
F(A) \;=\;
\frac{n}{2} \log R(A)
\;+\; \sum_{(i,j) \in \mathcal{S}} \log |A_{ij}|
\;+\; M \log (2 L_{A})
\;+\; \Pi(M)
\;-\; \log \Delta(R)
\end{array}}
\label{eq:criterion}
\end{equation}

where

\begin{equation}
\Delta(R) \;=\;
\Gamma\!\left( \frac{n}{2}, \frac{R}{2 s_{1}^{2}} \right)
- \Gamma\!\left( \frac{n}{2}, \frac{R}{2 s_{0}^{2}} \right)
\label{eq:Delta}
\end{equation}

The five terms have plain meanings. The first measures how badly the matrix fits the
data. The second and third concern the occupied coefficients: how large they are, and
how many. The fourth, given by \eqref{eq:Pi}, is the information required to specify which
cells are occupied.
The fifth is nearly constant over the ordinary range of $R$, by
\eqref{eq:sigmaordinary}, and takes effect only when the fit becomes so good as to be
implausible given the stated precision of the measurements.

Away from that extreme, $\Delta(R)$ may be replaced by $\Gamma(n/2)$ and dropped, so
that the working criterion is

\begin{equation}
F(A) \;=\;
\frac{n}{2} \log R(A)
\;+\; \sum_{(i,j) \in \mathcal{S}} \log |A_{ij}|
\;+\; M \log (2 L_{A})
\;+\; \Pi(M)
\label{eq:workingcriterion}
\end{equation}

One restriction on the use of \eqref{eq:criterion} must be stated at once. It is a
sound guide to the coefficient values within a pattern that has been fixed in
advance, and Section~\ref{sec:coefficients} uses it for that. It is not a sound guide
to the choice between patterns of different size, and Section~\ref{sec:compare}
explains why and gives the quantity that is.

\section{The coefficients for a given pattern}
\label{sec:coefficients}

Suppose the pattern $\mathcal{S}$ is fixed. The terms in
\eqref{eq:workingcriterion} that involve $M$ are then constant, and we minimise over
the values in the occupied cells. Write

\begin{equation}
e_{ij} \;=\; \sum_{t=1}^{T} \left( Y_{it} - Z_{it} \right) X_{jt}
\label{eq:eij}
\end{equation}

the inner product of the residuals in row $i$ with the $j$th input series, so that
$\partial R / \partial A_{ij} = -2 e_{ij}$. Setting the derivative of
\eqref{eq:workingcriterion} to zero gives, for every occupied cell,

\begin{equation}
e_{ij} \;=\; \frac{R(A)}{n \, A_{ij}}
\label{eq:stationary}
\end{equation}

Ordinary least squares would give $e_{ij} = 0$. The estimate here departs from least
squares by a quantity of order $R/n$, which is small when the measurements are many.
The direction of the departure is toward smaller coefficients, since the density
\eqref{eq:coefprior} places more of its value at small magnitudes.

That shrinkage deserves a word of caution. It is a consequence of maximising a
density rather than integrating one, and its size depends on the fact that we chose
to write the prior as a density in $A_{ij}$ rather than in $\log|A_{ij}|$. It is not
a statement about the data. Since it is of order $R/n$, it is negligible in the cases
of interest, and \eqref{eq:stationary} is retained as the estimate of the
coefficients throughout.

\section{Comparing patterns of different size}
\label{sec:compare}

Equation \eqref{eq:criterion} is not adequate for choosing between patterns, and the
reason is arithmetical. Occupying one more cell increases $\Pi$ by
\eqref{eq:increment}, which is under one natural unit in the middle of the range, and
increases $M \log(2 L_{A})$ by $\log (2 L_{A})$, which is about two. Against that,
one more coefficient reduces $R$, and the first term of \eqref{eq:criterion} carries
the factor $n/2$. When $T$ is small the reduction is substantial and the two charges
above do not approach it. Taken as a guide to the size of the pattern,
\eqref{eq:criterion} occupies as many cells as it is permitted to.

What is missing is any account of the freedom that the extra coefficient uses up.
Maximising over the coefficients records the height of the posterior at its best
point and disregards how wide the region of good fit is. A pattern with more
coefficients reaches a higher point, but over a region so much narrower that its
total weight may be smaller. The comparison must therefore be made between
integrals, not between heights.

\subsection{The integral over the coefficients}

Fix the pattern $\mathcal{S}$ and let $\hat{A}$ be the least-squares point within it,
$\hat{R}$ the residual sum of squares there. The residual is quadratic, so

\begin{equation}
R(A) \;=\; \hat{R} \;+\; d^{\,\prime} G \, d
\label{eq:quadratic}
\end{equation}

where $d$ collects the departures of the $M$ coefficients from $\hat{A}$ and $G$ is
block diagonal with blocks $G_{i} = X_{\mathcal{S}_{i}} X_{\mathcal{S}_{i}}^{\prime}$,
one for each row. The factor $\prod 1/|A_{ij}|$ varies slowly across the region where
the exponential is appreciable, so it may be evaluated at $\hat{A}$ and taken outside.
The remaining integral is Gaussian, and gives

\begin{equation}
\int \exp\!\left( -\frac{R(A)}{2\sigma^{2}} \right) dA
\;=\; \exp\!\left( -\frac{\hat{R}}{2\sigma^{2}} \right)
\left( 2\pi\sigma^{2} \right)^{M/2}
\prod_{i=1}^{N} \left| G_{i} \right|^{-1/2}
\label{eq:laplace}
\end{equation}

This requires each $G_{i}$ to be non-singular, which holds when no row has more than
$T$ occupied cells. A row with more than $T$ occupied cells reproduces its own data
exactly and leaves directions in which the residual does not change at all; such
patterns are excluded here.

\subsection{The integral over $\sigma$, once more}

The factor $\sigma^{M}$ from \eqref{eq:laplace} multiplies the $\sigma^{-(n+1)}$ of
Section~\ref{sec:sigmaint}, so the exponent becomes $M - n - 1$. Writing

\begin{equation}
\nu \;=\; \frac{n-M}{2}
\label{eq:nu}
\end{equation}

and repeating the substitution $u = 1/\sigma^{2}$,

\begin{equation}
\int_{s_{0}}^{s_{1}} \sigma^{\,M-n-1}
\exp\!\left( -\frac{\hat{R}}{2\sigma^{2}} \right) d\sigma
\;=\; \frac{1}{2} \left( \frac{2}{\hat{R}} \right)^{\nu}
\left[ \Gamma\!\left( \nu, \frac{\hat{R}}{2 s_{1}^{2}} \right)
     - \Gamma\!\left( \nu, \frac{\hat{R}}{2 s_{0}^{2}} \right) \right]
\label{eq:sigmaintM}
\end{equation}

Setting $M = 0$ recovers \eqref{eq:sigmaint} exactly.

\subsection{The criterion for the pattern}

Collecting the factors and taking the negative logarithm, the pattern is chosen by
minimising



\begin{equation}
\tcboxmath{ \begin{array}[t]{c}
\Phi(\mathcal{S}) \;=\;
\nu \log \hat{R}
\;+\; \sum_{(i,j) \in \mathcal{S}} \log |\hat{A}_{ij}|
\;+\; \tfrac{1}{2} \sum_{i=1}^{N} \log \left| G_{i} \right|
\;-\; \frac{M}{2} \log 2\pi
\;+\; M \log (2 L_{A})
\;+\; \Pi(M)
\end{array}}
\label{eq:patterncriterion}
\end{equation}


away from the extreme of a near-exact fit, where the bracket in \eqref{eq:sigmaintM}
must be retained in place of $\Gamma(\nu)$.

Two differences from \eqref{eq:workingcriterion} carry the whole effect. The factor
on $\log \hat{R}$ has fallen from $n/2$ to $(n-M)/2$, which is the count of
measurements less the count of quantities fitted to them. And two new terms appear,
$\tfrac{1}{2}\sum \log|G_{i}|$ and $-\tfrac{1}{2} M \log 2\pi$, which together measure
how narrow the region of good fit has become. A coefficient that the data determine
sharply makes $|G_{i}|$ large and raises $\Phi$; a coefficient that the data barely
constrain adds little. That is the accounting for freedom which \eqref{eq:criterion}
omits.

\subsection{A note on maximising against integrating}
\label{sec:integratenote}

The step taken in this section is narrow. Within a fixed pattern the coefficients are
still estimated by maximising, through \eqref{eq:stationary}; only the comparison
between patterns of different size is made by integration. It is worth recording that
the wider step may repay the effort.

The maximum of a density is a fragile quantity in two respects. It is not preserved
under a change of variable: writing the same distribution as a density in $\log|a|$
rather than in $a$ moves the maximum, while the integral over any region is
unaffected. And for a density that decreases monotonically, the maximum sits at an
endpoint whatever the parameters may be. The exponential density
$\lambda e^{-\lambda x}$ on $x \ge 0$ attains its maximum at $x = 0$ for every
$\lambda$, so the location of the maximum carries no information about $\lambda$ at
all. The prior \eqref{eq:coefprior} is of exactly this character, which is why
Section~\ref{sec:magnitudes} had to bound it and why \eqref{eq:stationary} shrinks the
coefficients slightly.

Wherever a conclusion rests on comparing regions of the parameter space rather than
on locating one point in it, the integral is the quantity with a stable meaning. The
present document integrates over $p$, over $\sigma$, and now over the coefficients
when patterns are compared. As the model is enlarged in later documents, results that
appear anomalous should be tested by integrating rather than maximising before any
other explanation is sought.

\section{Two numerical illustrations}
\label{sec:examples}

The algebra above says how a matrix is to be scored. Whether the score selects the
right matrix is a separate question, and this section answers it on manufactured data
where the truth is known.

Two matrices are used. In the first, the occupied cells form three blocks of size
three along the diagonal. In the second, the same coefficient values are moved to
three columns of each row chosen at random, so that no arrangement remains. Both
matrices therefore have the same number of occupied cells, the same values, and the
same count of three per row. They differ only in where those cells sit. The two cases
of a given trial share the same inputs and the same measurement errors, so the
comparison between them is paired.

Nothing in the prior refers to blocks. All $q = 81$ cells are treated alike, and the
rows are fitted separately. If the block structure is recovered, it is recovered from
the data.

\subsection{Construction}

\begin{center}
\begin{tabular}{ll}
dimension and observations & $N = 9$, \quad $T = 6$, \quad $q = 81$, \quad $n = 54$ \\
occupied cells & three per row, so $M = 27$ \\
coefficient magnitudes & drawn log-uniformly on $[0.40,\, 3.00]$, signs equally
likely \\
stated coefficient range & $a_{\min} = 0.20$, \quad $a_{\max} = 5.0$, \quad
$L_{A} = 3.219$ \\
measurement error & $\sigma = 0.10$ \\
stated error range & $s_{0} = 0.02$, \quad $s_{1} = 10$, \quad
$L_{\sigma} = 6.215$ \\
hyper-prior on $p$ & $\alpha = \beta = 1$, the uniform case \\
inputs & $X_{jt}$ drawn independently from the standard normal \\
\end{tabular}
\end{center}

The stated range $[a_{\min}, a_{\max}]$ lies strictly outside the range from which the
true coefficients are drawn. That separation matters. A floor placed at the edge of
the truth rejects those coefficients that the measurement error happens to push below
it, and the method then misses cells it would otherwise find.

The proportions of $N$ and $T$ are the point of the exercise. Each row supplies six
equations for nine candidate coefficients, so no row can be estimated without first
deciding which cells it occupies. Occupying three of them leaves three degrees of
freedom. Rows are restricted to at most $T-1 = 5$ occupied cells, which excludes exact
fits and keeps every $G_{i}$ of Section~\ref{sec:compare} non-singular.

\subsection{The search}

Section~\ref{sec:matrixspace} observed that a matrix of order $N$ admits $2^{N^{2}}$
patterns of occupied cells. For $N = 9$ that is $2^{81}$, or about
$2.4 \times 10^{24}$, so the patterns cannot be examined one by one and some
procedure is needed. This section states the one used.

Three properties of the problem make it tractable. The first is the restriction of
Section~\ref{sec:examples}.1 to at most $T-1 = 5$ occupied cells in a row, which is
required in any case for the residual to retain any degrees of freedom and for the
matrices $G_{i}$ of Section~\ref{sec:compare} to be non-singular. A row then admits

\begin{equation}
\sum_{k=0}^{5} \binom{9}{k} \;=\; 382
\label{eq:rowsubsets}
\end{equation}

subsets of columns rather than $2^{9} = 512$. The second is that the rows are fitted
separately, so the fit of a row depends on that row alone. Enumerating every
admissible subset of every row therefore takes $9 \times 382 = 3438$ least-squares
fits, and this is done once. For each row and each size $k$ the subset of that size
with the lowest cost is retained, which leaves six numbers per row.

The third property requires more care, because $\Phi$ does not separate across rows.
It couples them twice: through $\log \hat{R}$, which involves the sum of the residuals
of all nine rows, and through $\Pi(M)$, which involves the total count of occupied
cells. Holding $\sigma$ fixed removes the first coupling, since the joint posterior of
$(A,\sigma)$ contains $\hat{R}/2\sigma^{2}$ rather than its logarithm, and a sum
inside a logarithm becomes a sum of separate terms. What remains is a problem of
allocating a total of $M$ occupied cells among nine rows so as to minimise a sum of
per-row costs, and that is solved exactly by dynamic programming across the rows,
indexed by the number of cells allocated so far. One pass yields the best allocation
for every total $M$ at once.

The procedure is therefore: fix $\sigma$ at each of 240 values spaced logarithmically
across $[s_{0}, s_{1}]$; for each, run the dynamic programme and record the pattern it
gives for every $M$; discard duplicates; evaluate $\Phi$ exactly on every distinct
pattern that survives; and take the lowest. In the trial reported below this examined
862 distinct patterns in the block case and 790 in the scattered case. The coefficient
values of the winner are then moved from least squares to the stationary point of
\eqref{eq:stationary}.

What this does and does not guarantee should be stated plainly. For any fixed
$\sigma$ the dynamic programme returns the exact minimiser of the separable form, so
nothing is lost within a given $\sigma$. Across $\sigma$ the grid is finite, and
$\Phi$ itself is minimised only over the patterns the grid happens to produce. Some
$1.7 \times 10^{23}$ patterns satisfy the per-row restriction, of which a few hundred
are examined, so the procedure is a heuristic and not an exhaustive search.

That limitation cannot account for the behaviour reported below. A more complete
search would return a pattern of still lower $\Phi$, never a higher one, and every
departure from the truth reported here is a pattern whose $\Phi$ is already lower than
the truth attains. Searching harder would widen that gap rather than close it. The
over-fitting described in Section~\ref{sec:examples}.6 is thus a property of the
criterion and not an artefact of the procedure that minimises it.

\subsection{What can be compared with least squares, and what cannot}

For $Y = AX$ with $Y$ and $X$ of shape $N \times T$, minimising the sum of squared
residuals gives the normal equations $Y X^{\prime} = A X X^{\prime}$, hence

\begin{equation}
A \;=\; \left( Y X^{\prime} \right) \left( X X^{\prime} \right)^{-1}
\label{eq:leastsquares}
\end{equation}

Only one of the two possible comparisons can be made here. Since
$X X^{\prime} = \sum_{t} x_{t} x_{t}^{\prime}$ is a sum of $T$ matrices of rank one,
its rank is at most $T$ whatever the values in $X$. When $T \ge N$ and $X$ has full
row rank the inverse exists and \eqref{eq:leastsquares} is the estimate; a comparison
against it would then be meaningful. These illustrations use $N = 9$ and $T = 6$, so
the rank is at most six out of nine and the matrix is singular for every $X$, not
merely for unlucky ones. In the run reported below its three smallest singular values
are $1.7 \times 10^{-15}$, $9.8 \times 10^{-16}$ and $3.2 \times 10^{-16}$, which are
zeros carrying rounding error. The least-squares estimate does not exist, and no
comparison against it is possible.

What can be shown is the minimum-norm solution obtained from the pseudo-inverse. That
is a different estimate: among the matrices reproducing the data exactly it selects
the one of smallest Frobenius norm. 
Figure \ref{fig:LLSE_Matrix} show the linear least squares estimate of the matrix,
using the block structured matrix and data as given in section \ref{app:numbers}.
For the block case it occupies all 81 cells,
leaves a residual of $2 \times 10^{-28}$, and stands at Frobenius distance $4.02$ from
the true matrix.

\begin{figure}[ht]
\begin{center}
\footnotesize
\begin{tabular}{rrrrrrrrr}
  $-1.33$ & $1.21$ & $-0.64$ & $0.87$ & $-0.79$ & $0.61$ & $-0.31$ & $-0.48$ & $-0.44$ \\
  $-1.64$ & $0.95$ & $0.92$ & $0.25$ & $-0.25$ & $0.36$ & $-0.58$ & $-0.29$ & $-0.02$ \\
  $-1.78$ & $0.67$ & $0.86$ & $0.26$ & $-0.17$ & $0.23$ & $-0.30$ & $-0.25$ & $-0.05$ \\
  $0.38$ & $-0.21$ & $0.64$ & $-1.00$ & $1.32$ & $1.49$ & $0.35$ & $-0.12$ & $-0.47$ \\
  $-0.14$ & $-0.10$ & $-0.21$ & $1.18$ & $0.33$ & $-0.39$ & $0.08$ & $0.33$ & $0.54$ \\
  $0.16$ & $-0.33$ & $0.41$ & $-1.37$ & $-0.70$ & $-0.53$ & $-0.03$ & $-0.21$ & $-0.26$ \\
  $0.36$ & $0.35$ & $0.55$ & $0.32$ & $-0.04$ & $-0.35$ & $-0.96$ & $0.70$ & $0.34$ \\
  $-0.06$ & $-0.84$ & $-0.16$ & $0.26$ & $0.53$ & $-0.27$ & $0.42$ & $1.00$ & $-0.32$ \\
  $-0.12$ & $0.28$ & $-0.25$ & $-0.41$ & $-0.37$ & $0.42$ & $-0.14$ & $\cdot$ & $-0.64$ \\
\end{tabular}
\end{center}
% minimum-norm solution, block


\protect\caption{Linear Least Squares Estimated Matrix}
\label{fig:LLSE_Matrix}
\end{figure}


A comparison against \eqref{eq:leastsquares} itself requires a separate run with
$T \ge N$, which is not attempted here.

\subsection{One realisation}

The trial shown is the first of the forty, chosen before any results were seen. The
true block matrix and the matrix recovered from it are Figures
\ref{fig:True_Block_Matrix} and \ref{fig:Recovered_Block_Matrix}


\begin{figure}[ht]
\begin{center}
\footnotesize
\fbox{
\begin{tabular}{rrrrrrrrr}
  $-1.67$ & $2.30$ & $-1.30$ & $\cdot$ & $\cdot$ & $\cdot$ & $\cdot$ & $\cdot$ & $\cdot$ \\
  $-1.61$ & $1.72$ & $0.93$ & $\cdot$ & $\cdot$ & $\cdot$ & $\cdot$ & $\cdot$ & $\cdot$ \\
  $-1.75$ & $1.19$ & $0.77$ & $\cdot$ & $\cdot$ & $\cdot$ & $\cdot$ & $\cdot$ & $\cdot$ \\
  $\cdot$ & $\cdot$ & $\cdot$ & $-1.56$ & $1.82$ & $1.32$ & $\cdot$ & $\cdot$ & $\cdot$ \\
  $\cdot$ & $\cdot$ & $\cdot$ & $1.42$ & $0.55$ & $-0.72$ & $\cdot$ & $\cdot$ & $\cdot$ \\
  $\cdot$ & $\cdot$ & $\cdot$ & $-1.77$ & $-0.42$ & $-0.68$ & $\cdot$ & $\cdot$ & $\cdot$ \\
  $\cdot$ & $\cdot$ & $\cdot$ & $\cdot$ & $\cdot$ & $\cdot$ & $-1.44$ & $0.81$ & $0.98$ \\
  $\cdot$ & $\cdot$ & $\cdot$ & $\cdot$ & $\cdot$ & $\cdot$ & $1.63$ & $1.86$ & $0.44$ \\
  $\cdot$ & $\cdot$ & $\cdot$ & $\cdot$ & $\cdot$ & $\cdot$ & $-0.59$ & $-0.66$ & $-1.49$ \\
\end{tabular}
}
\end{center}
% true A, block


\protect\caption{True Block Matrix}
\label{fig:True_Block_Matrix}
\end{figure}


All 27 occupied cells are found, one further cell is occupied that should not be, and
the estimated error is $0.0798$ against a true $0.10$. The Frobenius distance from the
truth is $0.373$, against $4.02$ for the minimum-norm solution.

The blocks are visible in the recovered matrix, and nothing in the calculation was
told to look for them.

\begin{figure}[ht]
\begin{center}
\footnotesize
\begin{tabular}{rrrrrrrrr}
  $-1.65$ & $2.36$ & $-1.35$ & $\cdot$ & $\cdot$ & $\cdot$ & $\cdot$ & $\cdot$ & $\cdot$ \\
  $-1.59$ & $1.77$ & $0.88$ & $\cdot$ & $\cdot$ & $\cdot$ & $\cdot$ & $\cdot$ & $\cdot$ \\
  $-1.79$ & $1.22$ & $0.74$ & $\cdot$ & $\cdot$ & $\cdot$ & $\cdot$ & $\cdot$ & $\cdot$ \\
  $\cdot$ & $\cdot$ & $\cdot$ & $-1.57$ & $1.83$ & $1.31$ & $\cdot$ & $\cdot$ & $\cdot$ \\
  $\cdot$ & $\cdot$ & $\cdot$ & $1.41$ & $0.56$ & $-0.72$ & $\cdot$ & $\cdot$ & $\cdot$ \\
  $\cdot$ & $\cdot$ & $\cdot$ & $-1.76$ & $-0.38$ & $-0.71$ & $\cdot$ & $\cdot$ & $\cdot$ \\
  $\cdot$ & $-0.24$ & $\cdot$ & $\cdot$ & $\cdot$ & $\cdot$ & $-1.65$ & $0.72$ & $0.92$ \\
  $\cdot$ & $\cdot$ & $\cdot$ & $\cdot$ & $\cdot$ & $\cdot$ & $1.58$ & $1.88$ & $0.36$ \\
  $\cdot$ & $\cdot$ & $\cdot$ & $\cdot$ & $\cdot$ & $\cdot$ & $-0.60$ & $-0.69$ & $-1.49$ \\
\end{tabular}
\end{center}
% recovered A, block


\protect\caption{Recovered Block Matrix}
\label{fig:Recovered_Block_Matrix}
\end{figure}


For the scattered case the true and recovered matrices are in Figures
\ref{fig:True_Scattered_Matrix}
and
\ref{fig:Recovered_Scattered_Matrix}

\begin{figure}[ht]
\begin{center}
\footnotesize
\begin{tabular}{rrrrrrrrr}
  $\cdot$ & $\cdot$ & $-1.67$ & $\cdot$ & $\cdot$ & $\cdot$ & $2.30$ & $\cdot$ & $-1.30$ \\
  $\cdot$ & $-1.61$ & $1.72$ & $0.93$ & $\cdot$ & $\cdot$ & $\cdot$ & $\cdot$ & $\cdot$ \\
  $\cdot$ & $\cdot$ & $\cdot$ & $\cdot$ & $\cdot$ & $-1.75$ & $1.19$ & $\cdot$ & $0.77$ \\
  $\cdot$ & $-1.56$ & $\cdot$ & $1.82$ & $\cdot$ & $\cdot$ & $\cdot$ & $1.32$ & $\cdot$ \\
  $1.42$ & $\cdot$ & $0.55$ & $\cdot$ & $\cdot$ & $\cdot$ & $\cdot$ & $-0.72$ & $\cdot$ \\
  $-1.77$ & $\cdot$ & $\cdot$ & $\cdot$ & $\cdot$ & $\cdot$ & $-0.42$ & $-0.68$ & $\cdot$ \\
  $\cdot$ & $-1.44$ & $0.81$ & $\cdot$ & $\cdot$ & $\cdot$ & $\cdot$ & $0.98$ & $\cdot$ \\
  $\cdot$ & $\cdot$ & $\cdot$ & $\cdot$ & $1.63$ & $1.86$ & $\cdot$ & $0.44$ & $\cdot$ \\
  $\cdot$ & $\cdot$ & $\cdot$ & $\cdot$ & $-0.59$ & $-0.66$ & $-1.49$ & $\cdot$ & $\cdot$ \\
\end{tabular}
\end{center}
% true A, scatter


\caption{True Scattered Matrix}
\label{fig:True_Scattered_Matrix}
\end{figure}



Here the recovery is exact: all 27 cells found, none added, none missed, with an
estimated error of $0.0826$.

One number deserves attention. In the block case $\Phi$ at the recovered pattern is
$53.674$, while $\Phi$ at the true pattern is $55.919$. The criterion prefers the
pattern with one spurious cell to the truth, by more than two natural units. The
search did not fail; the criterion selected what it was asked to select.

\begin{figure}[ht]
\begin{center}
\footnotesize
\begin{tabular}{rrrrrrrrr}
  $\cdot$ & $\cdot$ & $-1.75$ & $\cdot$ & $\cdot$ & $\cdot$ & $2.26$ & $\cdot$ & $-1.27$ \\
  $\cdot$ & $-1.56$ & $1.67$ & $0.93$ & $\cdot$ & $\cdot$ & $\cdot$ & $\cdot$ & $\cdot$ \\
  $\cdot$ & $\cdot$ & $\cdot$ & $\cdot$ & $\cdot$ & $-1.72$ & $1.24$ & $\cdot$ & $0.82$ \\
  $\cdot$ & $-1.60$ & $\cdot$ & $1.81$ & $\cdot$ & $\cdot$ & $\cdot$ & $1.34$ & $\cdot$ \\
  $1.42$ & $\cdot$ & $0.53$ & $\cdot$ & $\cdot$ & $\cdot$ & $\cdot$ & $-0.71$ & $\cdot$ \\
  $-1.79$ & $\cdot$ & $\cdot$ & $\cdot$ & $\cdot$ & $\cdot$ & $-0.47$ & $-0.68$ & $\cdot$ \\
  $\cdot$ & $-1.51$ & $0.87$ & $\cdot$ & $\cdot$ & $\cdot$ & $\cdot$ & $0.98$ & $\cdot$ \\
  $\cdot$ & $\cdot$ & $\cdot$ & $\cdot$ & $1.60$ & $1.93$ & $\cdot$ & $0.50$ & $\cdot$ \\
  $\cdot$ & $\cdot$ & $\cdot$ & $\cdot$ & $-0.64$ & $-0.65$ & $-1.47$ & $\cdot$ & $\cdot$ \\
\end{tabular}
\end{center}
% recovered A, scatter


\protect\caption{Recovered Scattered Matrix}
\label{fig:Recovered_Scattered_Matrix}
\end{figure}

\subsection{Forty realisations}

One trial cannot distinguish a property of block structure from an accident of one
draw. The whole construction was therefore repeated over forty seeds, with the
coefficient values, the inputs, the errors and the scattered pattern all redrawn each
time.

\begin{center}
\begin{tabular}{lccccc}
case & correct of 27 & false & missed & exact & median $\hat\sigma$ \\
\hline
block & $26.20$ & $1.57$ & $0.80$ & $14$ of $40$ & $0.091$ \\
scatter & $26.40$ & $1.40$ & $0.60$ & $17$ of $40$ & $0.089$ \\
\end{tabular}
\end{center}
% paired: block worse 17, better 11, equal 12


The two cases perform alike. Compared within each seed, the block case was the worse
of the two seventeen times, the better eleven times, and equal twelve times; a sign
test on the twenty-eight untied pairs gives $p = 0.35$, so the difference is within
what chance produces.

That result is what the construction requires. The prior treats all cells alike and
the rows are fitted separately, so no information can pass from the arrangement of the
cells to the estimate. The method is indifferent to structure, which is what declining
to impose structure means. It also follows that the method does not exploit block
structure; it recovers such a matrix neither better nor worse than any other.

\subsection{What the illustrations show}

Three conclusions may be drawn, and a fourth qualification attaches to them.

The method recovers most of the correct cells. Across the forty trials it finds about
$26.3$ of the 27 occupied cells and misses fewer than one, in a setting where the
unrestricted estimate does not exist at all.

It recovers the whole matrix exactly in about two cases in five, and otherwise
occupies between one and two cells that should be empty. The estimated error is
consistently a little below the truth, near $0.09$ against $0.10$, which is the mark
of those extra cells absorbing part of the measurement error.

The residue is not distributed evenly. Repeating the forty trials with the floor
lowered to $a_{\min} = 0.05$ raises the number of falsely occupied cells from $1.49$
to $5.90$ per matrix, and the additional cells sit at small magnitudes. Their median
magnitude is $0.138$, against $1.171$ for the correctly occupied cells, and a third of
them fall below $0.10$, against fewer than one in two hundred of the correctly
occupied cells. Raising the floor to a value an analyst would assert removes most of
them and brings the estimated error back to $0.10$.

The qualification is that this concentration points at the approximation of
Section~\ref{sec:compare}. Equation \eqref{eq:laplace} treats
$\prod 1/|A_{ij}|$ as constant across the region of good fit and takes no account of
the truncation at $a_{\min}$. Both are least accurate for coefficients close to the
floor, and both err in the direction of overstating the integral there, which is
exactly where the spurious cells appear. Handling the truncation properly, rather than
by the Gaussian approximation, is the first improvement to attempt.

\section{Comparison with two established methods}
\label{sec:compete}

Section~\ref{sec:examples} compares the method with unrestricted least squares, which
does not exist at this size, and with the minimum-norm solution that stands in for
it. That solution occupies all eighty-one cells and is not a serious competitor. Two
methods that were designed for this regime are, and both were run on the same
illustrations. Full references for both are given in \texttt{literature-review.tex}
in this directory.

\subsection{The two methods}

The lasso, due to Tibshirani, adds to the residual sum of squares a penalty
proportional to the sum of the absolute values of the coefficients

\begin{equation}
\hat{a}(\lambda) \;=\; \arg\min_{a} \; \tfrac{1}{2} \| y - Z a \|^{2}
\;+\; \lambda \sum_{j=1}^{N} |a_{j}|
\label{eq:lasso}
\end{equation}

where $y$ is one row of $Y$ and $Z = X'$ is the $T$ by $N$ matrix of inputs. The
absolute value has a corner at the origin, and a penalty with a corner drives
coefficients to exactly zero rather than merely close to it, so that selection and
estimation are performed together and no search over patterns is required. The
minimiser of \eqref{eq:lasso} is the maximum posterior estimate under a
double-exponential prior on each coefficient, with $\lambda$ in the place of the
rate. It is therefore an estimate of the kind used here, differing in the prior
rather than in the principle. Two further differences matter. The prior is peaked at
the origin rather than bounded away from it, and there is no indicator of occupancy:
a cell is empty when the penalty happens to drive it to zero, and not because a
choice was made about it. The penalty itself is not supplied by the model and must be
chosen by other means.

Stochastic search variable selection, due to George and McCulloch and carried to
vector autoregressions by George, Sun and Ni, attaches an indicator $\gamma_{ij}$ to
each cell and makes the prior on the coefficient depend upon it

\begin{eqnarray}
A_{ij} \mid \gamma_{ij} = 0 &\sim& N(0, \tau_{0}^{2})
\label{eq:ssvs_A}
\\
A_{ij} \mid \gamma_{ij} = 1 &\sim& N(0, \tau_{1}^{2})
\label{eq:ssvs_B}
\\
\gamma_{ij} &\sim& \mathrm{Bernoulli}(p)
\label{eq:ssvs_C}
\end{eqnarray}

with $\tau_{0}$ small and $\tau_{1}$ large. Neither component is exactly zero, so
every coefficient always has a value, every conditional distribution in the model is
a standard one, and the posterior can be explored by drawing in turn the
coefficients, the indicators, the rate $p$ and the error variance. No move that
changes the number of parameters is ever required, which is what makes the sampler
practical. The output is a sample of patterns rather than a single pattern, and the
summary taken here is the set of cells occupied in at least half the draws.

The resemblance to the present work is close. The pattern prior is the same Bernoulli
form of Section~\ref{sec:patterncost}, the hyper-prior on $p$ is the same Beta, one
error scale serves the whole matrix, and the hierarchy is arranged in the same order.
Two things differ. The coefficient prior is a pair of normal densities rather than
the bounded log-uniform of Section~\ref{sec:magnitudes}, and the pattern is sampled
rather than maximised.

\subsection{What was compared}

Both methods were run on the two illustrations over the same forty seeds, with the
same matrices, the same inputs and the same measurement errors, so that the
comparison is paired at the level of the seed and at the level of the case. That
gives eighty paired comparisons. The script is \texttt{examples/compete.py}.

Four variants of the lasso were run, because one variant would not settle the matter.
The first chooses the penalty by leave-one-out cross-validation, which is what an
analyst without an estimate of the noise would do. The second chooses it by the
Schwarz criterion, given the true error scale, which no analyst would have. The third
is the adaptive lasso of Zou, reweighted from the second, which penalises each cell
in inverse proportion to its magnitude in that fit. The fourth chooses the penalty,
row by row, by counting the pattern errors against the truth and taking the penalty
that makes them fewest; this requires the answer and is therefore not a method but a
bound on what any choice of penalty could achieve. All four have their coefficients
refitted by least squares on the pattern they select, so that they are charged for
the pattern rather than for the shrinkage.

Stochastic search variable selection was run with $\tau_{0} = 0.02$ and
$\tau_{1} = 1.50$, twenty thousand draws of which the first five thousand were
discarded. Those widths stand in the same relation to the true magnitudes, drawn
log-uniformly on $[0.40, 3.00]$, as the bounds $[0.20, 5.0]$ used here, so that
neither method is given the better information.

\subsection{Forty realisations, six procedures}

\begin{center}
\footnotesize
\begin{tabular}{llccccc}
case & method & correct of 27 & false & missed & exact & median $\hat\sigma$ \\
\hline
block & this paper & $26.20$ & $1.57$ & $0.80$ & $14$ of $40$ & $0.091$ \\
block & lasso, cross-validated & $20.85$ & $18.43$ & $6.15$ & $0$ of $40$ & $0.614$ \\
block & lasso, Schwarz, true $\sigma$ & $23.70$ & $18.27$ & $3.30$ & $0$ of $40$ & $0.065$ \\
block & adaptive lasso, Schwarz & $23.32$ & $8.28$ & $3.67$ & $0$ of $40$ & $0.092$ \\
block & lasso, penalty from the answer & $19.27$ & $1.68$ & $7.72$ & $0$ of $40$ & $0.746$ \\
block & SSVS & $26.32$ & $1.50$ & $0.68$ & $10$ of $40$ & $0.091$ \\
\hline
scatter & this paper & $26.40$ & $1.40$ & $0.60$ & $17$ of $40$ & $0.089$ \\
scatter & lasso, cross-validated & $21.73$ & $19.30$ & $5.28$ & $0$ of $40$ & $0.430$ \\
scatter & lasso, Schwarz, true $\sigma$ & $23.98$ & $17.00$ & $3.02$ & $0$ of $40$ & $0.065$ \\
scatter & adaptive lasso, Schwarz & $23.57$ & $7.85$ & $3.42$ & $2$ of $40$ & $0.096$ \\
scatter & lasso, penalty from the answer & $19.68$ & $2.35$ & $7.33$ & $0$ of $40$ & $0.696$ \\
scatter & SSVS & $26.00$ & $1.12$ & $1.00$ & $17$ of $40$ & $0.094$ \\
\end{tabular}
\end{center}


The first line of each half reproduces the forty-realisation table of
Section~\ref{sec:examples} exactly, which confirms that the procedure compared
against is the one described in this document.

Counting a pattern error as a cell wrongly occupied or a cell wrongly empty, the
eighty paired comparisons fall out as follows.

\begin{center}
\footnotesize
\begin{tabular}{lcccc}
method & worse & better & equal & sign test \\
\hline
lasso, cross-validated & $80$ & $0$ & $0$ & $2 \times 10^{-24}$ \\
lasso, Schwarz, true $\sigma$ & $80$ & $0$ & $0$ & $2 \times 10^{-24}$ \\
adaptive lasso, Schwarz & $77$ & $2$ & $1$ & $1 \times 10^{-20}$ \\
lasso, penalty from the answer & $79$ & $0$ & $1$ & $3 \times 10^{-24}$ \\
SSVS & $24$ & $24$ & $32$ & $1.00$ \\
\end{tabular}
\end{center}
% 80 paired comparisons, errors counted as false plus missed


\subsection{What the comparison shows}

No variant of the lasso is competitive here, and the variant given the answer is not
competitive either. The best of the four occupies more than eight cells that should
be empty and misses nearly four that should be occupied, against a little over two
errors in total for the method of this document.

The reason is not obscure. With six observations and nine candidates in each row, a
single continuous penalty must serve two purposes at once: it must be large enough to
empty six cells and small enough to retain three. The fourth variant shows that no
choice of penalty does both, since even with the answer in hand it must accept nearly
eight missed cells in order to hold its false count down. Cross-validation, which is
aimed at prediction rather than at recovering the pattern, fails in the opposite
direction and occupies about nineteen cells too many.

Against stochastic search variable selection the result is even. The method of this
document wins twenty-four of the eighty comparisons, loses twenty-four, and ties
thirty-two, and the sign test gives $p = 1.00$. The average number of pattern errors
is $2.19$ here against $2.15$ for the sampler. The estimate of the error scale is
$0.090$ against $0.093$, with the truth at $0.100$. The criterion developed here
attains the accuracy of the established method at this size, and does not exceed it.

Two differences are nevertheless real. The first is the worst case. The sampler goes
wrong less often, exceeding five pattern errors in six of the eighty comparisons
against eleven here, but when it does go wrong it goes further: its worst trial
misplaces twenty-one cells, and the worst trial here misplaces ten. The second is
that the procedure here
is deterministic. It evaluates a criterion in closed form on candidates produced by a
dynamic programme, returns the same answer whenever it is run, and can be read term
by term. That last property is what made the account in Section~\ref{sec:examples} of
where the residue comes from possible at all, and a sample of patterns yields no such
account.

One further point should be recorded, since it bears on the objection that
$a_{\min}$ and $a_{\max}$ are not determined by the data. The two widths
$\tau_{0}$ and $\tau_{1}$ are not determined by the data either, and the sampler is
at least as sensitive to them. Repeating the first illustration at four settings,
with two chains at each, the pair $(0.05, 1.00)$ and the pair $(0.02, 1.50)$ recover
the matrix, while $(0.005, 3.00)$ leaves two cells occupied out of twenty-seven and
$(0.10, 0.60)$ occupies all eighty-one. The failures run in both directions, and the
first of them is the divergence of Section~\ref{sec:magnitudes} appearing inside the
sampler: spreading the prior over a wider range lowers the marginal likelihood of
every occupied cell. The script is \texttt{examples/ssvs\_sensitivity.py}.

This does not dispose of the objection. It establishes that the requirement is a
property of the problem rather than of the formulation adopted here, and that the two
bounds have an advantage of form, being stated in the units of the coefficients as
the smallest and largest magnitude worth calling non-zero, which is a quantity an
analyst can be asked to supply.

\section{The search at two hundred variables}
\label{sec:large}

Section~\ref{sec:examples} used nine variables and enumerated every admissible
subset of the columns for every row. This section reports what happens at the problem size
which  an economist might encounter. We imagine a hypothetical
data set for twenty countries of ten variables each,
so that $A$ is $200$ by $200$ with twenty ten-by-ten diagonal blocks.
For sizing, we assumed that each country is connected by two variables
to every other country (a different pair from each influencing country,
which is its connection to the other nineteen). We assumed
twenty years of quarterly data.
The complete matrix would forty thousand cells, of which we expect about $2500-2800$ to be non-zero.
Again, we have a prior on how many are non-zero, but (as noted earlier) the methodology
does not even have a way to represent block structure, much less require it.
The matrix is to be estimated from eighty quarters of data. Each quarter is
represented simply by a $200$-vector at the end of the quarter which is estimated
from the $200$-vector at the start of the quarter. As noted earlier, we do not use the time-structure in any way.
Enumeration is not merely slow at that size; it would require more than
the age of the universe. The section says why, describes what replaced it,
and reports the result.

\subsection{Why enumeration cannot be used}

The search of Section~\ref{sec:examples} lists every subset of the columns of size
at most $T-1$, that limit being what keeps $G_{i}$ non-singular. At $N = 200$ and
$T = 80$ the number of such subsets is about $3 \times 10^{57}$ for each row, so the
list is never built and the procedure never reaches the data.

Lowering the limit does not help, because the limit has to be at least as large as
the number of cells a row actually holds. Here that number is about thirteen, and
the count of subsets of size at most thirteen drawn from two hundred columns is
$9.5 \times 10^{19}$. Fitting all of them for all two hundred rows is
$2 \times 10^{22}$ least-squares problems. At the measured cost of $28$
microseconds each that is $5 \times 10^{17}$ seconds, or seventeen billion
years, which is longer than the age of the universe.

The limits that are affordable are far too small. A limit of two admits twenty
thousand subsets for each row and runs in under a minute. A limit of three admits
$1.3 \times 10^{6}$, which would run in about forty minutes but needs some thirty
gigabytes to hold the coefficients of every subset for every row. A limit of four is
out of reach on both counts. Three cells to a row is less than a quarter of what the
matrix carries.

It is worth being clear about which part of the procedure fails. The dynamic
programme is not the difficulty. With the limit set at twenty-five, its arrays are
of length $M_{\max} = 5000$ and it costs a few seconds over the whole grid of
$\sigma$. Evaluating $\Phi$, refining the coefficients and forming the residuals are
all matters of seconds. The obstacle is the list of subsets that the dynamic
programme selects from, and that list alone.

\subsection{Forward selection}

The dynamic programme needs one thing from each row: for every size $k$, a subset of
that size, with its residual sum of squares, its sum of $\log |A_{ij}|$ and its
$\log |G_{i}|$. Enumeration supplies the best subset of each size. Forward selection
supplies a subset of each size at a cost that grows as $N$ rather than as
$\binom{N}{k}$.

The rule is the ordinary one. Beginning with no cells, admit at each step the column
that reduces the residual sum of squares by the most, and stop at $K$ cells. Holding
an orthonormal basis $Q$ of the columns already admitted makes the reduction
available for every candidate at once. Writing $r$ for the current residual and
$Z_{j}$ for a candidate column, the reduction obtained by admitting $j$ is

\begin{equation}
\Delta_{j} \;=\; \frac{( Z_{j}' r )^{2}}
                      {\| Z_{j} \|^{2} - \| Q' Z_{j} \|^{2}}
\label{eq:forward}
\end{equation}

so that one matrix product settles the whole step. This is the exact forward step
and not an approximation to it. The work for a row is $K$ such steps, and the whole
matrix of two hundred rows is prepared in about one second.

What is given up should be stated plainly. The subset of size $k$ is now the one
that forward selection reaches, not the best subset of size $k$. The criterion is
therefore minimised over a chain of $K$ nested subsets for each row rather than over
all subsets, and the answer is a lower bound on how well $\Phi$ could be made to do.
Given those candidate subsets the dynamic programme is exact, as before, so the
approximation is confined to one place and can be named.

Two remarks temper the loss. The enumeration of Section~\ref{sec:examples} was
already limited to subsets of size at most $T-1$, so it was never a search over all
patterns either. And forward selection is not an arbitrary choice: it is the
standard method for this problem, and the coefficients here are well separated from
zero, which is the case in which it does well.

\subsection{Two further changes forced by the size}

The integral over $\sigma$ is now taken in the closed form \eqref{eq:sigmaint}
rather than by quadrature. The integrand of \eqref{eq:sigmaintdef}, written in
$u = \log \sigma$, is concentrated about its peak in a range of width
$1 / \sqrt{2(n-M)}$. At $n = 54$ that width is about $0.14$, and the grid used in
Section~\ref{sec:examples} put some forty points across it. At $n = 16000$ the width
is about $0.006$ and the same grid puts four. The closed form has no such
difficulty, costs less, and was verified in Maxima along with everything else.

The criterion is also accumulated through the dynamic programme instead of being
recomputed on patterns that have been built. Inspection of
\eqref{eq:workingcriterion} shows that $\Phi$ depends on a pattern only through $M$
and three quantities that are sums over the rows: the residual sum of squares, the
sum of $\log |A_{ij}|$ and the sum of $\log |G_{i}|$. Carrying those three sums
alongside the cost lets $\Phi$ be evaluated for every pair of $\sigma$ and $M$
without any pattern being formed, and only the winner is turned into a matrix. At
this size the alternative would be to build and store some $1.2$ million patterns of
two hundred entries each.

\subsection{The problem}

Each of the twenty countries is a solid block of ten by ten coefficients, which is
two thousand cells. Each ordered pair of countries is joined by one or two cells
drawn at random from the hundred that connect them, which adds about five hundred
and seventy. The matrix therefore holds about $2570$ occupied cells out of forty
thousand, or one cell in fifteen, and a row holds about thirteen. The occupied
fraction is close to $2/29$, so the hyper-prior of Section~\ref{sec:patterncost} is
taken as $\alpha = 2$ and $\beta = 27$, whose mean is $2/29$ and whose density is
proportional to $p (1-p)^{26}$. As noted earlier, the prior does not force a specific answer
and the actual number estimated is determined by the quality of the fit and amount of
data, using the Occam criterion.

Everything else is as in Section~\ref{sec:examples}. The magnitudes are drawn
log-uniformly on $[0.40, 3.00]$, the bounds are $[0.20, 5.0]$, the error scale is
$0.10$, and $K$ is set at twenty-five. A scattered control is generated by keeping
the number of cells in each row and placing them at random columns, so that the two
matrices differ in arrangement alone: one is $200 \times 200$ with twenty ten-by-ten diagonal
blocks, and the other is  $200 \times 200$ with no block structure at all. We will show
that the MPLE method works equally well on both.

One point deserves emphasis before the numbers. Per row this problem is easier than
the illustration of Section~\ref{sec:examples}, not harder: eighty observations
carry about thirteen coefficients, a ratio of six, against six observations carrying
three, a ratio of two. The difficulty at this size lies wholly in the search.

\subsection{Results}

The procedure of this document was run over ten seeds, in both arrangements. Each
matrix holds about $2569$ occupied cells.

\begin{center}
\footnotesize
\begin{tabular}{llccccccc}
case & method & trials & occupied & correct of & false & missed & median $\hat\sigma$ & median error \\
\hline
block & this paper & $10$ & $2551$ & $2541$ of $2569$ & $10$ & $28$ & $0.263$ & $3.96$ \\
\hline
scatter & this paper & $10$ & $2552$ & $2544$ of $2569$ & $8$ & $25$ & $0.232$ & $3.92$ \\
\end{tabular}
\end{center}


It recovers about $99$ per cent of the occupied cells. The two arrangements are
again alike: compared within the seed the block case was the worse of the two seven
times and the better three times, which a sign test does not distinguish from
chance. That is the result of Section~\ref{sec:examples} once more, and for the same
reason. Nothing in the criterion or in the search can see the arrangement.

\subsection{Where the missing cells go}
\label{sec:whymissed}

The cells that are not recovered are worth following, because they turn out to be
the price of the substitution made in this section and not a fault of the criterion.

The first fact is that the criterion is not the difficulty. Evaluating
\eqref{eq:workingcriterion} at the true pattern, with the coefficients fitted by
least squares, gives $\Phi = -6355$, while the pattern the procedure returns gives
$\Phi = 5093$. The criterion prefers the truth by more than eleven thousand natural
units, an enormous margin, and almost all of it comes through the residual sum of
squares, which is $137$ at the true pattern and $767$ at the recovered one. Whatever
went wrong, the criterion would have rejected it had it been offered the alternative.

The second fact is that forward selection is very nearly good enough. For $194$ of
the $200$ rows the subset it reaches at that row's own size is exactly the true
subset of that row. The limit $K$ is not binding either, the fullest row of the truth
holding eighteen cells against a limit of twenty-five.

The third fact explains the first two. In the six rows where forward selection admits
a wrong column, the least-squares coefficient of that column is small, and it falls
below $a_{\min}$. Section~\ref{sec:magnitudes} places no density below $a_{\min}$, so
the subset containing it is not merely poor but inadmissible, and the dynamic
programme cannot use that size for that row at all. It must fall back on a smaller
size, which discards good cells along with the bad one. Six rows truncated in this
way carry about six hundred of the excess residual, which is the whole of the
difference between $137$ and $767$, and through
Section~\ref{sec:sigmaint} that excess is what raises $\hat\sigma$ from $0.10$ to
about $0.25$.

The mechanism is therefore a single ordering error in a few rows, turned into a large
loss by the rule that rejects a subset outright when one of its coefficients falls
outside the range of the prior. Note that the true subset of a row is not what is
rejected: it passes the test in all two hundred rows. What is rejected is the subset
forward selection offers in place of it.

That is a defect of the procedure adopted here, and it is repaired by a better search
and not by any change to the criterion. Section~\ref{sec:candidates} repairs it.

\subsection{The comparison at this size}
\label{sec:largecompare}

The six procedures of Section~\ref{sec:compete} were run again at this size, on the
same data, with the settings there except that the cross-validation is in five folds
rather than one for each observation, which at $T = 80$ is the usual practice.

\begin{center}
\footnotesize
\begin{tabular}{llccccccc}
case & method & trials & occupied & correct of & false & missed & median $\hat\sigma$ & median error \\
\hline
block & this paper & $6$ & $2552$ & $2541$ of $2570$ & $11$ & $29$ & $0.256$ & $3.79$ \\
block & lasso, cross-validated & $6$ & $11441$ & $2570$ of $2570$ & $8870$ & $0$ & $0.048$ & $2.21$ \\
block & lasso, Schwarz, true $\sigma$ & $6$ & $5340$ & $2570$ of $2570$ & $2770$ & $0$ & $0.084$ & $1.49$ \\
block & adaptive lasso, Schwarz & $6$ & $2570$ & $2570$ of $2570$ & $0$ & $0$ & $0.100$ & $0.63$ \\
block & lasso, penalty from the answer & $6$ & $2590$ & $2252$ of $2570$ & $338$ & $318$ & $0.829$ & $14.30$ \\
block & SSVS & $6$ & $4060$ & $2270$ of $2570$ & $1790$ & $301$ & $0.177$ & $43.53$ \\
\hline
scatter & this paper & $6$ & $2553$ & $2544$ of $2570$ & $9$ & $26$ & $0.247$ & $3.96$ \\
scatter & lasso, cross-validated & $6$ & $11467$ & $2570$ of $2570$ & $8896$ & $0$ & $0.048$ & $2.24$ \\
scatter & lasso, Schwarz, true $\sigma$ & $6$ & $5366$ & $2570$ of $2570$ & $2796$ & $0$ & $0.084$ & $1.53$ \\
scatter & adaptive lasso, Schwarz & $6$ & $2570$ & $2570$ of $2570$ & $0$ & $0$ & $0.100$ & $0.63$ \\
scatter & lasso, penalty from the answer & $6$ & $2599$ & $2258$ of $2570$ & $341$ & $312$ & $0.796$ & $13.84$ \\
scatter & SSVS & $6$ & $4042$ & $2278$ of $2570$ & $1764$ & $292$ & $0.177$ & $41.92$ \\
\end{tabular}
\end{center}


The result is the reverse of Section~\ref{sec:compete} and it must be reported as
such. The adaptive lasso recovers the matrix exactly. In every trial, in both
arrangements, it occupies precisely the cells that are occupied in the truth, misses
none, adds none, and returns $\hat\sigma = 0.10$ against a truth of $0.10$. The
procedure of this document recovers about $99$ per cent of the cells and misses two
dozen. At nine variables no variant of the lasso was competitive; at two hundred, one
of them is exact and this criterion is not.

Two things should be said about that before it is read as a verdict on the criterion.

The first is that it is not one. Section~\ref{sec:whymissed} shows that
$\Phi$ at the true pattern is smaller than $\Phi$ at the pattern recovered here by
more than eleven thousand natural units. The criterion ranks the answer the adaptive
lasso found far above the answer this search found. What the comparison establishes
is that the adaptive lasso is the better search at this size, not that it is the
better criterion. The natural thing to try next is to use the lasso path as the
source of candidate patterns and $\Phi$ to choose among them, which is not done here.

The second is a qualification that has been tested and does not hold. The two lasso
variants scored by the Schwarz criterion were given the true error scale, which no
analyst would have, and it was natural to suspect that the exactness was bought with
that. It was not. Estimating the error scale by the ordinary plug-in iteration --
choosing the penalty of the plain lasso by cross-validation to obtain a first value,
then alternating between choosing the pattern at the current value and re-estimating
the value from the residuals of the chosen pattern -- reaches the same answer.

\begin{center}
\footnotesize
\begin{tabular}{llrrccc}
seed & case & starting $\hat\sigma$ & $M$ at the start & passes & correct & final $\hat\sigma$ \\
\hline
$20260726$ & block & $0.0462$ & $11563$ & $1$ & $2573$ of $2573$ & $0.1009$ \\
$20260726$ & block & $1.0000$ & --- & $2$ & $2573$ of $2573$ & $0.1009$ \\
$20260726$ & scatter & $0.0456$ & $11601$ & $1$ & $2573$ of $2573$ & $0.1010$ \\
$20261726$ & block & $0.0472$ & $11351$ & $1$ & $2553$ of $2553$ & $0.0996$ \\
$20261726$ & scatter & $0.0492$ & $11299$ & $1$ & $2553$ of $2553$ & $0.0998$ \\
\end{tabular}
\end{center}
% no falsely occupied cell and no missed cell in any row of this table;
% true error scale 0.10; from examples/large_alasso_sigma.txt


The iteration begins about four times too dense, at some eleven thousand cells with
an error scale less than half the truth, and one pass carries it to the exact
pattern. Started instead from a deliberately large value of unity it arrives at the
same place in two. The fixed point does not depend on where the iteration begins, and
it is the truth. Whatever else is to be said about this comparison, the adaptive
lasso is not being flattered by information an analyst would lack.

The plain lasso behaves as it did at the smaller size, in one respect. Every variant
of it finds all the occupied cells; what none of them does, except the adaptive one,
is stop. Cross-validation occupies some eleven and a half thousand cells against a
truth of about two and a half thousand, and the Schwarz criterion about five
thousand. The variant given the answer, which chooses the penalty row by row to make
the fewest errors, is markedly worse than the adaptive variant, which shows that the
advantage lies in the reweighting rather than in any information about the truth.

Stochastic search variable selection cannot be assessed from this table. The figures
shown are what the sampler reached in eight thousand sweeps, and
Section~\ref{sec:mixing} shows that at this size it is nowhere near its own answer at
that length or at twenty thousand. They are a statement about a fixed amount of
computing and not about the method.

\subsection{Whether the sampler was run long enough}
\label{sec:mixing}

It was not. The sampler was run at four lengths on the first trial, twice at each
from different random numbers.

\begin{center}
\footnotesize
\begin{tabular}{rrcccccc}
draws & burn & chain & $M$ & correct & false & missed & chains differ on \\
\hline
$4000$ & $1000$ & 1 & $4820$ & $2135$ & $2685$ & $438$ & \\
$4000$ & $1000$ & 2 & $4725$ & $2159$ & $2566$ & $414$ & $5503$ cells \\
$8000$ & $2000$ & 1 & $4481$ & $2193$ & $2288$ & $380$ & \\
$8000$ & $2000$ & 2 & $4165$ & $2250$ & $1915$ & $323$ & $4424$ cells \\
$20000$ & $5000$ & 1 & $4019$ & $2274$ & $1745$ & $299$ & \\
$20000$ & $5000$ & 2 & $3891$ & $2290$ & $1601$ & $283$ & $3518$ cells \\
$60000$ & $15000$ & 1 & $3647$ & $2348$ & $1299$ & $225$ & \\
$60000$ & $15000$ & 2 & $3602$ & $2342$ & $1260$ & $231$ & $2721$ cells \\
\end{tabular}
\end{center}
% seed 20260726, block, true cells 2573; from examples/large_ssvs_mix.txt
% and examples/large_ssvs_mix_60000.txt


Every quantity is still moving toward the truth as the run lengthens, and the two
chains still disagree about several thousand of the forty thousand cells, with
individual inclusion probabilities differing by almost one. That is not a sampler
settled about an answer; it is one still travelling toward it.

Worse, it is slowing down. Counting the falsely occupied cells in excess of the
truth, each doubling of the run removes about $450$ of them between four and eight
thousand sweeps, about $280$ between eight and twenty thousand, and about $210$
between twenty and sixty thousand. Fitting a power of the run length to those four
points gives an exponent near $-0.27$, on which the excess would fall to a few dozen
cells only after some thousands of millions of sweeps. That extrapolation should not
be taken literally, since four points are a thin basis for judging the convergence of
a Markov chain, but the direction of it is not in doubt. At sixty thousand sweeps,
three times the length of the longest published application of this method and on a
problem sixty-seven times larger, the sampler still occupies about $1280$ cells that
should be empty.

The comparison in Section~\ref{sec:compete} at nine variables gave this method and
the sampler an exact tie. At two hundred variables the sampler cannot be run long
enough, within any time that was available here, to say what it would conclude. That
is a real difference between a criterion evaluated in closed form and a distribution
explored by sampling, and it is the sharpest form of the point made in
Section~\ref{sec:compete}. It should not, however, be read as a defect of stochastic
search variable selection as a method. The largest published application of it that
could be found carries $595$ coefficients, and the sampler used here is the plain
one; blocked and scalable variants exist and none was tried.

\subsection{Where the candidates come from}
\label{sec:candidates}

Section~\ref{sec:whymissed} places the difficulty in the search and not in the
criterion, and Section~\ref{sec:largecompare} shows that the adaptive lasso finds the
pattern this search does not. That suggests keeping $\Phi$ as the judge and changing
only the list of patterns it is asked to judge. It was tried, and it answers the
question.

The dynamic programme requires, for each row and each size, a subset with its
residual sum of squares, its sum of $\log |A_{ij}|$ and its $\log |G_{i}|$. Forward
selection supplies one such subset for each size. A penalty path supplies others: it
is indexed by penalty rather than by size, so it visits some sizes several times and
others not at all, but every support it visits is a candidate. Nothing prevents the
two lists being merged, and since $\Phi$ decides among them, a longer list cannot
give a worse answer. The programme was therefore given several candidates of each
size and takes at each $\sigma$ whichever the row cost prefers, which is what the
enumeration of Section~\ref{sec:examples} did at nine variables.

\begin{center}
\footnotesize
\begin{tabular}{llccccccc}
case & candidates from & trials & occupied & correct of & false & missed & median $\hat\sigma$ & median $\Phi$ \\
\hline
block & forward selection & $3$ & $2554$ & $2547$ of $2569$ & $7$ & $22$ & $0.239$ & --- \\
block & and the plain path & $3$ & $2555$ & $2547$ of $2569$ & $8$ & $22$ & $0.239$ & $5095$ \\
block & and the adaptive path & $3$ & $2569$ & $2569$ of $2569$ & $0$ & $0$ & $0.101$ & $-6398$ \\
\hline
scatter & forward selection & $3$ & $2552$ & $2545$ of $2569$ & $7$ & $24$ & $0.270$ & --- \\
scatter & and the plain path & $3$ & $2553$ & $2546$ of $2569$ & $7$ & $23$ & $0.269$ & $6642$ \\
scatter & and the adaptive path & $3$ & $2569$ & $2569$ of $2569$ & $0$ & $0$ & $0.100$ & $-6431$ \\
\end{tabular}
\end{center}
% three seeds, both arrangements; from examples/large_lasso_phi.txt
% no falsely occupied and no missed cell in any of the six adaptive runs


The plain path adds almost nothing. The reason is that it seldom passes through the
right subset: over the two hundred rows of the first trial it visits the true subset
of a row nineteen times, against one hundred and ninety-four for forward selection.

The adaptive path recovers the matrix exactly. In all six trials, in both
arrangements, every occupied cell is found, no empty cell is occupied, and
$\hat\sigma$ returns to $0.10$. The value of $\Phi$ falls from about $5100$ to about
$-6400$, and that second figure is the value of $\Phi$ at the true pattern reported
in Section~\ref{sec:whymissed}. The criterion has not merely improved; it has found
the pattern it said all along that it preferred.

One point of construction decides whether this works. The weights of the adaptive
path must come from a fit that is too dense, not from one that is too sparse.
Weighting by the answer $\Phi$ has just given fails completely, because a cell that
answer leaves empty receives the largest weight the rule allows and can never be
recovered. Weighting instead by a cross-validated fit, which occupies some eleven
thousand cells and therefore contains every cell that ought to be occupied, gives
each true cell a small weight and each false one a large weight. Cross-validation
requires no estimate of the error scale, and it is used here only to weight and never
to select, so no quantity enters that an analyst would lack. Every selection is made
by $\Phi$.

The cost is the one drawback. Forward selection prepares a matrix in about a second
and the whole procedure runs in sixteen; the cross-validated weighting takes some
forty minutes, being five penalty paths for each of two hundred rows. That is worth
reducing and there is no obstacle of principle to reducing it, but it is the present
position.

\subsection{The direction of the error}

The error has changed direction with the size of the problem. In
Section~\ref{sec:examples}, with $q = 81$ and $M$ near $27$, the procedure occupied
one or two cells more than it should. Here, with $q = 40000$ and $M$ near $2570$, it
leaves about twenty-six fewer than it should.

Part of this is arithmetic that would hold whatever search were used. The cost of one
more occupied cell is $\log \left[ (q-M-1+\beta) / (M+\alpha) \right]$, which is
$0.64$ natural units at the smaller size and $2.68$ at the larger. Adding
$\log 2 L_{A} = 1.86$, a cell must return about $4.5$ natural units before it is
worth occupying here, against about $2.5$ before, so the criterion is more demanding
in a large matrix than in a small one.

That is not, however, what produced the omissions reported above. Section
\ref{sec:whymissed} shows that the criterion prefers the true pattern by a very wide
margin and would have taken it. The omissions here are the work of the search. The
two effects should not be confused: the toll per cell does rise with $q$, and the
cells missed in this study were nevertheless missed for a different reason.

What is not available is a proof. Nothing in this document establishes whether the
criterion occupies more cells than the truth or fewer, on average, nor under what
conditions the one turns into the other. Two sizes have been examined and they fall
on opposite sides, and an arithmetic account has been given of why the toll per cell
rises with $q$; that account is consistent with the observations but it is not a
theorem, and it says nothing about the size of the discrepancy. Whether
$E[\hat{M}] - M$ is positive, negative or zero in general, how it behaves as $T$
grows with $q$ fixed, and whether the pattern is recovered in the limit, are all
open. They are the natural questions to take up next, and they are questions about
the criterion rather than about the search.

\subsection{The severity of the test}

The count used throughout is a strict one, and the numbers should be read with that
in mind. A cell is scored wrong if it is occupied and should be empty, or empty and
should be occupied, and no allowance is made for magnitude. Missing a coefficient of
$0.4$ counts exactly as much as missing one of $3.0$. No credit is given for
recovering the arrangement, for placing the twenty blocks correctly, or for coming
close. The test asks for the exact set of forty thousand decisions and records every
departure.

By that measure the block case is charged with thirty-eight errors in a matrix of
$2569$ occupied cells. By almost any other measure the same result would be
described differently. The distance between the estimate and the truth has a median
of $3.96$ over the whole matrix, which is about $0.078$ for each occupied cell
against a typical magnitude of $1.1$. Most of that distance is the missed cells
themselves: twenty-six omitted coefficients of median magnitude $0.625$ account for
some three units of the four, so the coefficients that were retained are estimated
very closely indeed. A user interested in forecasting, or in which countries
influence which, would find little of the error visible.

The strict count is nevertheless the right one to report, because the question this
document asks is which cells are occupied, and a test of that question should not
be softened by allowing partial credit for nearly answering it. But it is a severe
test, and it is not the test by which such models are usually judged.

\section{What is not settled here}

Three matters are deliberately left open.

The search over patterns is not addressed. There are $2^{q}$ of them, and $q$ may be
forty thousand. Equation \eqref{eq:workingcriterion} says how to score a pattern; it
does not say how to find a good one. Note that, before $p$ and $\sigma$ are removed,
the $N$ rows of $A$ are separate problems, each involving only its own row of $Y$.
The rows are connected solely through the two shared quantities. That separation is
the natural starting point for a search procedure, and it is the whole of the
hierarchical structure in the present model.

The second point is that one $p$ and one $\sigma$ serve the entire matrix. In the
AutoClass formalism, parameters are shared within classes rather than across
everything. Allowing $p$ and $\sigma$ to vary by row, or by group of rows, is the
natural next step and will be taken up separately.

Third, the inputs are taken to be exact. Error in $X$ as well as $Y$ is a different
and harder problem and is not considered.

\section{Verification}

Every algebraic step above was checked with Maxima. The script is
\texttt{maxima/acv\_prelim.mac} in this directory, and it verifies:

\begin{enumerate}
\item the divergence of $\int da/|a|$ at both ends, equations
      \eqref{eq:diverge_A} and \eqref{eq:diverge_B};
\item that the bounded densities \eqref{eq:coefprior} and \eqref{eq:sigmaprior_B}
      integrate to unity;
\item that $L_{A}$ is unchanged by a change of units;
\item the integral over $p$, equation \eqref{eq:pint}, for general $\alpha$ and
      $\beta$; its reduction \eqref{eq:pintuniform} at $\alpha = \beta = 1$, in all
      three of the forms written there, symbolically and at particular values of $q$
      and $M$; the ratio \eqref{eq:ratio} and hence the increment
      \eqref{eq:increment}; and the posterior mean of $p$;
\item the closed form \eqref{eq:sigmaint} for the integral over $\sigma$, checked
      against numerical quadrature at three sets of values;
\item the finite limit \eqref{eq:sigmalimit} as $R \to 0$, and the divergence of the
      corresponding unbounded expression;
\item the location \eqref{eq:peak} of the maximum of the integrand in
      \eqref{eq:sigmanaive}, the unbounded growth of its height as $R \to 0$, and the
      closed form \eqref{eq:sigmaunbounded}, the last checked against numerical
      quadrature over $(0,\infty)$ at $R = 1,\ 0.1,\ 0.01,\ 0.001$;
\item the stationarity condition \eqref{eq:stationary};
\item the Gaussian integral behind \eqref{eq:laplace}, and the closed form
      \eqref{eq:sigmaintM} for the integral over $\sigma$ with the exponent
      $M-n-1$, checked against numerical quadrature at three values of $\nu$ and
      confirmed to reduce to \eqref{eq:sigmaint} at $M = 0$.
\end{enumerate}

The illustrations of Section~\ref{sec:examples} are produced by the scripts in the
\texttt{examples} directory. \texttt{recover.py} carries the first trial together
with the comparison between the criterion of Section~\ref{sec:criterion} and that of
Section~\ref{sec:compare}; \texttt{sweep.py} repeats the construction over many
seeds; \texttt{final.py} produces the numbers and the tables printed above. The
search is implemented twice, independently, and the two implementations agree on the
scattered case of the first seed. That agreement is the check each was written
against, and it is what exposed an error in the first: it collected candidate
patterns under a key that recorded only how many cells each row received, so that
only the first value of $\sigma$ producing a given allocation contributed a pattern,
and part of the candidate set went unexamined.

\appendix
% The report class attaches \appendix to chapters.  This document numbers by
% section, so the section counter is reset and lettered here instead.
\setcounter{section}{0}
\renewcommand{\thesection}{\Alph{section}}

\pagebreak{}


\section{Numerical Examples}
\label{app:numbers}

Every number needed to repeat the calculations of Section~\ref{sec:examples} is
printed here: the two matrices to be recovered, the inputs, and the outputs of both
cases. A reader who works from these tables alone reaches the same results, and
needs neither our programs nor the generator that produced the figures.

The arrays were nonetheless produced from a single stream of pseudo-random numbers,
started at seed $20260726$ and drawn in the order given below, using the default
generator of the \texttt{numpy} library, version 2.5. The script
\texttt{examples/final.py} carries out all of it. That recipe is recorded for anyone
wishing to extend the example to other draws; it is not needed to check the one
reported.

\begin{enumerate}
\item Twenty-seven magnitudes are drawn, each with $\log$ uniform on
      $[\log 0.40,\; \log 3.00]$.
\item Twenty-seven signs are drawn, positive and negative equally likely.
\item Each magnitude is multiplied by its sign and rounded to two decimal places.
      These twenty-seven values are the occupied cells of both matrices.
\item The values are placed row by row into three blocks of size three along the
      diagonal, giving the first matrix.
\item For the second matrix, each row keeps its own three values and moves them to
      three columns drawn without replacement from the nine. A row that happens to
      draw its own block columns causes the whole arrangement to be drawn again.
\item The nine by six array of inputs $X$ is drawn from the standard normal.
\item A nine by six array of measurement errors is drawn from the standard normal
      and multiplied by $\sigma = 0.10$.
\end{enumerate}

The outputs are then $Y = AX + E$, formed separately for each of the two matrices from
the same $X$ and the same $E$.

\subsection{The block matrix}

The occupied cells lie in three blocks of size three along the diagonal. Empty cells
are shown as a dot.

\begin{center}
\footnotesize
\fbox{
\begin{tabular}{rrrrrrrrr}
  $-1.67$ & $2.30$ & $-1.30$ & $\cdot$ & $\cdot$ & $\cdot$ & $\cdot$ & $\cdot$ & $\cdot$ \\
  $-1.61$ & $1.72$ & $0.93$ & $\cdot$ & $\cdot$ & $\cdot$ & $\cdot$ & $\cdot$ & $\cdot$ \\
  $-1.75$ & $1.19$ & $0.77$ & $\cdot$ & $\cdot$ & $\cdot$ & $\cdot$ & $\cdot$ & $\cdot$ \\
  $\cdot$ & $\cdot$ & $\cdot$ & $-1.56$ & $1.82$ & $1.32$ & $\cdot$ & $\cdot$ & $\cdot$ \\
  $\cdot$ & $\cdot$ & $\cdot$ & $1.42$ & $0.55$ & $-0.72$ & $\cdot$ & $\cdot$ & $\cdot$ \\
  $\cdot$ & $\cdot$ & $\cdot$ & $-1.77$ & $-0.42$ & $-0.68$ & $\cdot$ & $\cdot$ & $\cdot$ \\
  $\cdot$ & $\cdot$ & $\cdot$ & $\cdot$ & $\cdot$ & $\cdot$ & $-1.44$ & $0.81$ & $0.98$ \\
  $\cdot$ & $\cdot$ & $\cdot$ & $\cdot$ & $\cdot$ & $\cdot$ & $1.63$ & $1.86$ & $0.44$ \\
  $\cdot$ & $\cdot$ & $\cdot$ & $\cdot$ & $\cdot$ & $\cdot$ & $-0.59$ & $-0.66$ & $-1.49$ \\
\end{tabular}
}
\end{center}
% true A, block



\subsection{The scattered matrix}

The same twenty-seven values, each row keeping its own three, moved to columns drawn
at random.

\begin{center}
\footnotesize
\begin{tabular}{rrrrrrrrr}
  $\cdot$ & $\cdot$ & $-1.67$ & $\cdot$ & $\cdot$ & $\cdot$ & $2.30$ & $\cdot$ & $-1.30$ \\
  $\cdot$ & $-1.61$ & $1.72$ & $0.93$ & $\cdot$ & $\cdot$ & $\cdot$ & $\cdot$ & $\cdot$ \\
  $\cdot$ & $\cdot$ & $\cdot$ & $\cdot$ & $\cdot$ & $-1.75$ & $1.19$ & $\cdot$ & $0.77$ \\
  $\cdot$ & $-1.56$ & $\cdot$ & $1.82$ & $\cdot$ & $\cdot$ & $\cdot$ & $1.32$ & $\cdot$ \\
  $1.42$ & $\cdot$ & $0.55$ & $\cdot$ & $\cdot$ & $\cdot$ & $\cdot$ & $-0.72$ & $\cdot$ \\
  $-1.77$ & $\cdot$ & $\cdot$ & $\cdot$ & $\cdot$ & $\cdot$ & $-0.42$ & $-0.68$ & $\cdot$ \\
  $\cdot$ & $-1.44$ & $0.81$ & $\cdot$ & $\cdot$ & $\cdot$ & $\cdot$ & $0.98$ & $\cdot$ \\
  $\cdot$ & $\cdot$ & $\cdot$ & $\cdot$ & $1.63$ & $1.86$ & $\cdot$ & $0.44$ & $\cdot$ \\
  $\cdot$ & $\cdot$ & $\cdot$ & $\cdot$ & $-0.59$ & $-0.66$ & $-1.49$ & $\cdot$ & $\cdot$ \\
\end{tabular}
\end{center}
% true A, scatter



\subsection{The inputs}

The nine by six array $X$. Both cases use it unchanged, so it is given once. The row
index $j$ runs over the components of the input vector and the column index $t$ over
the six observations.

\begin{center}
\footnotesize
\fbox{
\begin{tabular}{r|rrrrrr}
   & $t=1$ & $t=2$ & $t=3$ & $t=4$ & $t=5$ & $t=6$ \\
\hline
  $j=1$ & $-0.166450$ & $-0.721720$ & $-1.956128$ & $0.498841$ & $-1.404545$ & $-0.285642$ \\
  $j=2$ & $0.516260$ & $0.647030$ & $-0.876677$ & $-0.527945$ & $0.386305$ & $1.130293$ \\
  $j=3$ & $-1.164318$ & $1.145200$ & $0.142888$ & $-1.157197$ & $-0.554896$ & $-1.233712$ \\
  $j=4$ & $1.711768$ & $-0.854384$ & $-0.920882$ & $0.134728$ & $2.363427$ & $-0.971539$ \\
  $j=5$ & $-1.785136$ & $-0.735074$ & $-0.423137$ & $-0.314053$ & $1.523255$ & $-1.244201$ \\
  $j=6$ & $-1.925274$ & $-0.040422$ & $-0.242413$ & $-0.164816$ & $1.433408$ & $1.335336$ \\
  $j=7$ & $-0.580635$ & $-1.296970$ & $0.915671$ & $-0.097830$ & $0.186884$ & $0.458883$ \\
  $j=8$ & $0.030004$ & $0.803816$ & $0.392157$ & $1.644819$ & $1.237441$ & $-2.244700$ \\
  $j=9$ & $0.844496$ & $-0.704454$ & $-1.322737$ & $-1.225774$ & $-0.544313$ & $-0.788596$ \\
\end{tabular}
}
\end{center}
% inputs X, shared by both cases



\subsection{The outputs}

The nine by six array $Y$ for the block case, formed as $Y = AX + E$ with $A$ the
matrix of Section~\ref{app:numbers}.1 and $E$ the measurement errors. The row index
$i$ runs over the components of the output vector.

\begin{center}
\footnotesize
\fbox{
\begin{tabular}{r|rrrrrr}
   & $t=1$ & $t=2$ & $t=3$ & $t=4$ & $t=5$ & $t=6$ \\
\hline
  $i=1$ & $3.089032$ & $1.152337$ & $0.945229$ & $-0.521665$ & $4.004424$ & $4.778189$ \\
  $i=2$ & $0.160786$ & $3.342429$ & $1.686812$ & $-2.720215$ & $2.428925$ & $1.354611$ \\
  $i=3$ & $0.066173$ & $2.819019$ & $2.494934$ & $-2.510754$ & $2.645237$ & $0.923847$ \\
  $i=4$ & $-8.470004$ & $-0.074040$ & $0.441555$ & $-0.879926$ & $0.977153$ & $0.986482$ \\
  $i=5$ & $2.800136$ & $-1.651273$ & $-1.298047$ & $0.206664$ & $3.134637$ & $-3.015192$ \\
  $i=6$ & $-0.948718$ & $1.893286$ & $1.935759$ & $-0.126420$ & $-5.763637$ & $1.223792$ \\
  $i=7$ & $1.646148$ & $1.947265$ & $-2.188468$ & $0.302910$ & $-0.001640$ & $-3.394931$ \\
  $i=8$ & $-0.701104$ & $-0.797361$ & $1.623413$ & $2.503228$ & $2.384561$ & $-3.805208$ \\
  $i=9$ & $-0.854054$ & $1.217295$ & $1.098340$ & $0.871725$ & $-0.169908$ & $2.499949$ \\
\end{tabular}
}
\end{center}
% outputs Y, block case



The same for the scattered case, using the matrix of
Section~\ref{app:numbers}.2 and the same $E$:

\begin{center}
\footnotesize
\fbox{
\begin{tabular}{r|rrrrrr}
   & $t=1$ & $t=2$ & $t=3$ & $t=4$ & $t=5$ & $t=6$ \\
\hline
  $i=1$ & $-0.378846$ & $-4.032070$ & $3.467584$ & $3.322333$ & $2.113083$ & $4.238572$ \\
  $i=2$ & $-1.154212$ & $0.135980$ & $0.713241$ & $-0.947912$ & $0.640831$ & $-4.747311$ \\
  $i=3$ & $3.385596$ & $-2.110847$ & $0.500295$ & $-0.890321$ & $-2.550343$ & $-2.369103$ \\
  $i=4$ & $2.340323$ & $-1.518996$ & $0.304329$ & $3.359343$ & $5.331900$ & $-6.521777$ \\
  $i=5$ & $-0.933285$ & $-1.036591$ & $-2.913672$ & $-1.042962$ & $-3.217771$ & $0.542177$ \\
  $i=6$ & $0.540249$ & $1.320387$ & $2.774337$ & $-2.092267$ & $1.600207$ & $1.728883$ \\
  $i=7$ & $-1.698984$ & $0.902531$ & $1.871200$ & $1.465823$ & $0.005519$ & $-4.969843$ \\
  $i=8$ & $-6.659635$ & $-0.788115$ & $-0.984590$ & $0.047925$ & $5.711317$ & $-0.563054$ \\
  $i=9$ & $3.270531$ & $2.325823$ & $-1.028175$ & $0.513020$ & $-2.177188$ & $-0.716799$ \\
\end{tabular}
}
\end{center}
% outputs Y, scattered case



The errors themselves are not printed. They are the difference between these outputs
and $AX$, and nothing in the calculation uses them: the method sees only $X$ and $Y$.

\subsection{The precision of the printed values}

The coefficients of $A$ are exact as printed, since the generator rounds them to two
decimal places. The entries of $X$ and $Y$ are given to six decimal places, so a
reader recomputing from them works with slightly different numbers from those used
above.

That rounding does not affect the result. Repeating the whole calculation on the
rounded arrays recovers the same set of occupied cells in both cases, and the
estimated coefficients differ from those reported by less than $10^{-6}$. Six decimal
places are therefore sufficient, and the check is carried out by
\texttt{examples/final.py} rather than asserted.

\end{document}
